% concatenated from Domes_arXiv_v2.tex
% for section-numbered lemmas etc., use "numberwithinsect"
%\documentclass[a4paper,english,numberwithinsect,thm-restate]{eurocg24-submission}
%\documentclass[a4paper,english,numberwithinsect,thm-restate]{eurocg24}
%eurocg24.cls
\documentclass[11pt]{article}
\pdfoutput=1 %Unclear if needed.

% the recommended bibstyle
%\bibliographystyle{plainurl}
\bibliographystyle{plain}

\usepackage{hyperref}
\usepackage{url,float}
\usepackage{graphicx}
\usepackage{amsmath}
\usepackage{amsfonts}
\usepackage{amssymb}
\usepackage{latexsym}
%\usepackage{mathtools}
\usepackage{color}
\usepackage{xcolor}
%\usepackage[dvipsnames]{xcolor}
%\usepackage{yfonts}
%\usepackage[boxruled]{algorithm2e}
\usepackage[shortlabels]{enumitem}
%\begin{enumerate}[(a)]

%\usepackage{subfigure}
%\usepackage{caption}
%\usepackage{subcaption}

\usepackage{hyperref}
\hypersetup{
    colorlinks=true,
    linkcolor=blue,
    filecolor=magenta,      
    urlcolor=cyan,
    }
%\usepackage{appendix}

\newcommand{\hide}[1]{}

\newcommand{\ABox}{
\raisebox{3pt}{\framebox[6pt]{\rule{6pt}{0pt}}}
}
% Anna: I see that this (below) defines a new proof environment. I'd favour using the default, but at least I added a \medkip at the end of the proof, as is standard.
\newenvironment{proof}{{\bf Proof:}}{\hfill\ABox\medskip}
%\newenvironment{proof}{{\bf Proof:}}{\hfill\ABox}

\newtheorem{theorem}{{\bf Theorem}}
\newtheorem{corollary}[theorem]{Corollary}
\newtheorem{proposition}{Proposition}
\newtheorem{lemma}[theorem]{Lemma}
\newtheorem{conjecture}[theorem]{Conjecture}
\newtheorem{definition}[theorem]{Definition}
\newtheorem{obs}[theorem]{Observation}
\newtheorem{claim}[theorem]{Claim}

%For labels of items.
\newcommand{\lemlab}[1]{\label{lemma:#1}}
\newcommand{\thmlab}[1]{\label{thm:#1}}
\newcommand{\exlab}[1]{\label{ex:#1}}
\newcommand{\proplab}[1]{\label{prop:#1}}
\newcommand{\corlab}[1]{\label{cor:#1}}
%\newcommand{\deflab}[1]{\label{def:#1}}
\newcommand{\tablab}[1]{\label{tab:#1}}
\newcommand{\figlab}[1]{\label{fig:#1}}
\newcommand{\seclab}[1]{\label{sec:#1}}
\newcommand{\rmklab}[1]{\label{rmk:#1}}
\newcommand{\eqnlab}[1]{\label{eqn:#1}}

\newcommand{\lemref}[1]{\ref{lemma:#1}}
\newcommand{\thmref}[1]{\ref{thm:#1}}
\newcommand{\corref}[1]{\ref{cor:#1}}
\newcommand{\propref}[1]{\ref{prop:#1}}
\newcommand{\exref}[1]{\ref{ex:#1}}
%\newcommand{\defref}[1]{\ref{def:#1}}
%\newcommand{\chapref}[1]{\ref{chap:#1}}
\newcommand{\secref}[1]{\ref{sec:#1}}
\newcommand{\eqnref}[1]{\ref{eqn:#1}}
\newcommand{\figref}[1]{\ref{fig:#1}}
\newcommand{\tabref}[1]{\ref{tab:#1}}
\newcommand{\rmkref}[1]{\ref{rmk:#1}}


\def\a{{\alpha}}
\def\b{{\beta}}
\def\q{{\theta}}
% Blackboard R for real numbers, S for sphere (\S taken someplace, so \Sph).
\def\P{{\mathcal P}}
%\def\r{{\rho}}
%\def\s{{\sigma}}
%\def\S{{\Sigma}}
% \def\Q{{\mathcal Q}}
\def\D{{\mathcal D}}
% \def\X{{\mathcal X}}
%\def\G{{\Gamma}}
%\def\g{{\gamma}}
%\def\l{{\lambda}}
% \def\L{{\Lambda}}
% \def\k{{\kappa}}
\def\o{{\omega}}
%\def\O{{\Omega}}
\def\d{{\delta}}
\def\e{{\varepsilon}}
%\def\I{{\iota}}
%\def\hull{\mathop{\rm hull}\nolimits}
%\def\sp{\mathop{\rm sp}\nolimits}
%Special, problematical symbols:
%\def\polyh{{\mathcal P}}
%\def\p{{P}}
%\def\p*{{P^*}}%\def\p1{{P_1}}
%\def\p2{{P_2}}
% \def\bP{{\partial P}}
% \def\bX{{\partial X}}
% \def\bD{{\partial D}}
% \def\bU{{\partial U}}
% \def\bM{{\partial M}}
\def\R{{\mathbb{R}}}
\def\S{{\mathbb{S}}}
%\def\Sk{{\operatorname {Sk}}}

\def\defn#1{\textit{\textbf{\boldmath #1}}}

\newcommand{\squeezelist}{\setlength{\itemsep}{0pt}}
%\newcounter{abc}

%Commenting; suggested text
\newcommand{\JOR}[1]{{\color{red} JOR: [#1]}}
\newcommand{\jor}[1]{{\color{cyan} #1}}
\newcommand{\Anna}[1]{{\color{red} Anna: [#1]}}
\newcommand{\anna}[1]{{\color{blue} #1}}
\newcommand{\pepa}[1]{{\color{brown} #1}}
\newcommand{\old}[1]{{\color{gray} Old: [#1]}}

%Camera-ready: black, or {}
%\newcommand{\JOR}[1]{{}}
%\newcommand{\jor}[1]{{\color{black} #1}}
%\newcommand{\Anna}[1]{{}}
%\newcommand{\anna}[1]{{\color{black} #1}}
%\newcommand{\pepa}[1]{{\color{black} #1}}

% Support an insane number of symbol footnotes, for extreme numbers of author
% affiliations.  This code supports 6x3 = 18 symbol footnotes, using the 6
% standard symbols repeated between once and thrice.  Standard LaTeX stops
% around half way through that sequence.
{\makeatletter
\gdef\@fnsymbol#1{\ensuremath{\ifcase#1\or *\or \dagger\or \ddagger\or
   \mathsection\or \mathparagraph\or \|\or **\or \dagger\dagger\or
   \ddagger\ddagger\or \mathsection\mathsection\or
   \mathparagraph\mathparagraph\or \|\|\or ***\or \dagger\dagger\dagger\or
   \ddagger\ddagger\ddagger\or \mathsection\mathsection\mathsection\or
   \mathparagraph\mathparagraph\mathparagraph\or \|\|\|\else\@ctrerr\fi}}}


\title{%
Deltahedral Domes over Equiangular Polygons\footnote{%
Full version of abstract presented at EuroCG~\cite{DomesEuroCG}.
J.T. was supported by the Center for Foundations of Modern Computer Science (Charles Univ. project UNCE/SCI/004) and by project PRIMUS/24/SCI/012 from Charles Univ.}
} %title

%\author[1]{Antonio Example}
%\author[2]{Sylvia Veryknown}
%\affil[1]{Important University\\
%  \texttt{aexample@cs.important.de}}
%\affil[2]{Department of Mathematics and Computer
%        Science, University of Somewhere\\
%  \texttt{sveryknown@math.somewhere.nl}}
%%mandatory. First: Use abbreviated first/middle names.
%%Second (only in severe cases): Use first author plus 'et al.'
%\authorrunning{A. Example and S. Veryknown}

%\author[1]{MIT CompGeom Group}
%\affil[1]{Artificial 1st author to highlight that %the other authors
%%(in alphabetical order) 
%coauthors worked as an equal group.}
%%Please include all authors (including this one) in your bibliography,
%%and refer to the authors as ``MIT CompGeom Group'' (without ``et al.'').}
%%
%\author[2]{Hugo A. Akitaya}
%\affil[2]{U. Mass. Lowell, \texttt{hugo\_akitaya@uml.edu}}
%%
%\author[3]{Erik D. Demaine}
%\affil[3]{MIT, \texttt{edemaine@mit.edu}}
%%
%\author[4]{Adam Hesterberg}
%\affil[4]{Harvard U., \texttt{achesterberg@gmail.com}}
%%
%\author[5]{Anna Lubiw}
%\affil[5]{U. Waterloo, \texttt{alubiw@uwaterloo.ca}}
%%
%\author[6]{Jayson Lynch}
%\affil[6]{MIT, \texttt{jaysonl@mit.edu}}
%%
%\author[7]{Joseph O'Rourke}
%\affil[7]{Smith College, \texttt{jorourke@smith.edu}.}
%%
%\author[8]{Frederick Stock}
%\affil[8]{U. Mass. Lowell, \texttt{fbs9594@rit.edu}}
%%
%\author[9]{Josef Tkadlec}
%\affil[9]{Charles U., \texttt{josef.tkadlec@iuuk.mff.cuni.cz}}
%%
%\authorrunning{MIT CompGeom Group}
%Alpha-order:
\author{%
MIT CompGeom Group\thanks{%
Artificial first author to highlight that the other authors
(in alphabetical order) worked as an equal group.
Please include all authors (including this one) in your bibliography,
and refer to the authors as ``MIT CompGeom Group'' (without ``et al.'').}
\and
Hugo A. Akitaya\thanks{%
U. Mass. Lowell, \texttt{hugo\_akitaya@uml.edu}}
\and
Erik D. Demaine\thanks{%
MIT, \texttt{edemaine@mit.edu}}
\and
Adam Hesterberg\thanks{%
Harvard U., \texttt{achesterberg@gmail.com}}
\and
Anna Lubiw\thanks{%
U. Waterloo, \texttt{alubiw@uwaterloo.ca}}
\and
Jayson Lynch\thanks{%
MIT, \texttt{jaysonl@mit.edu}}
\and
Joseph O'Rourke\thanks{%
Smith College, \texttt{jorourke@smith.edu}}
\and
Frederick Stock\thanks{%
U. Mass. Lowell, \texttt{fbs9594@rit.edu}} 
\and
Josef Tkadlec\thanks{%
Charles U.
\texttt{jtkadlec@ist.ac.at}}
}%author


\begin{document}

\date{}
\maketitle



\begin{abstract}
A \emph{polyiamond} is a polygon composed of unit equilateral triangles, and
a \emph{generalized deltahedron}
is a convex polyhedron %where 
whose every face is a convex polyiamond.
We study a variant where one face may be an exception.
For a convex polygon $P$, 
if there is a convex polyhedron that has $P$ as one face and all the other faces are convex polyiamonds, then we say that 
% A \emph{delta dome} is a convex polyhedron where every face but one is a convex polyiamond.  If $P$ is the exceptional face, we say that
$P$ \emph{can be domed}.
Our main result is a complete characterization of 
which equiangular $n$-gons can be domed:
only if $n \in \{3,4,5,6,8,10,12\}$, and only with some 
% "certain" -> "some" to get down to a 4 line abstract
conditions on the integer edge lengths.
\end{abstract}



\section{Introduction}
\seclab{Introduction}
In the study of what can be built with equilateral triangles, the most well-known result is that there are exactly eight convex \defn{deltahedra}---polyhedra where every face is an equilateral triangle---with $n = 4,5,6,7,8,9,10,12$ vertices.
See references in \cite{bezdek2007proof} or
Wikipedia.\footnote{\url{https://en.wikipedia.org/wiki/Deltahedron}}
% \Anna{give reference, maybe wikipedia
% \url{https://en.wikipedia.org/wiki/Deltahedron}, or see refs. in [2]}.
What if coplanar triangles are allowed?
%We will allow coplanar equilateral triangles. 
In the plane, the polygons built of equilateral triangles are the \defn{polyiamonds}. Convex polyiamonds have 3, 4, 5, or 6 vertices.
The convex polyhedra with polyiamond faces are the 
``non-strictly convex deltahedra'', or \defn{generalized deltahedra}, following the nomenclature of Bezdek~\cite{bezdek2007proof}.
See the above cited Wikipedia article
for %some 
examples. 
%\Anna{How about (to keep it brief, but hopefully accurate): 
There are an infinite number of generalized deltahedra, though  the number of combinatorial types is finite since they have at most $12$ vertices.
%Although there are an infinite number of generalized deltahedra, they have at most $12$ vertices, so the number of combinatorial types is finite.
% Generalized deltahedra have
% at most $12$ vertices, so the number of combinatorial types is finite, though the class is infinite.
% \JOR{Not clear what is a ``class.'' My try:}
% \jor{though there are an infinite number of incongruent
% generalized deltahedra.}
There is no published characterization, though a forthcoming one is mentioned in~\cite{bezdek2007proof}.




% Our goal is to characterize when a convex polygon can be ``domed'' with a convex surface composed of equilateral triangles.  More formally, a \defn{delta dome} is a convex polyhedron where every face but one is a convex polyiamond. If $P$ is the exceptional face, we say that $P$ \defn{can be domed}, and that $P$ is the \defn{base} of the dome. Our goal is to characterize the bases of delta domes.
% \JOR{I think this def is not accurate, because we
% also dome polyiamonds. Here's my try:}
% \jor{...a \defn{delta dome} is a convex polyhedron
% $P \cup \D$
% with a convex polygon \defn{base} face $P$, and
% all other faces convex polyiamonds,
% forming the dome $\D$.}
% \Anna{Oh, I see the issue! 
% The dome should just be the non-base part.  The abstract needs fixing too.  We actually never use ``delta dome'' except to define ``doming''.
% How about:
Our goal (only partially achieved) 
is to characterize when a convex polygon can be ``domed'' with a convex surface composed of equilateral triangles.  
For a convex polygon $P$, 
if there is a convex polyhedron that has $P$ as one face and all the other faces are convex polyiamonds, then we say that $P$ \defn{can be deltahedrally domed}, or just \defn{domed} for short.  
%\Anna{If we use the title ``Deltahedral Domes over Equiangular Polygons'' then we could say that $P$ can be deltahedrally domed and that $\D$ is the deltahedral dome (dome for short).}
Here the \defn{deltahedral dome} (\defn{dome} for short), denoted by $\D$, is the part of the polyhedron excluding face $P$, and $P$ is called the \defn{base} of the dome.  Note that $P$ itself may or may not be a polyiamond.



%\hide{
%\subsection{Delta Domes}
%Let $P$ be a convex polygon lying in the $xy$-plane.
%We define a \emph{delta dome $\D$ over $P$} 
%(or just a \emph{dome} for short)
%to be a convex polyhedral surface composed of unit equilateral triangles 
%such that $P \cup \D$ is a convex polyhedron. We say that then $P$ can be \emph{domed}.
%We require that none of the triangles of $\D$ lie in the $xy$-plane:
%all have at least one vertex in the open halfspace $z > 0$.
%So $P \cap \D = \partial P$, and the polyhedron has positive volume.
%}

We assume that all the equilateral triangles have unit edge length,
so $P$ must be an integer polygon (with integer side lengths).
%
Here is a simple example:
\begin{lemma}
\lemlab{Rect}
Every integer rectangle can be domed.
\end{lemma}
%
\begin{proof}
If the base $P$ is an $a \times b$ rectangle, $a \ge b$,
then one can construct a ``roof" dome
whose top ridge has length $a-b$, and whose opposite
faces are trapezoids and triangles.
See Figs.~\figref{RectInt} and~\figref{Rect_3x1_roof}.
\end{proof}


%%%%%%%%%%%%%%%%%%%%%%%%%%%%%%%%%Figure Begin
\begin{figure}[htbp]
\centering
\includegraphics[width=0.55\textwidth]{Figures/RectInt}
%\includegraphics[width=0.75\textwidth]{Figures/RectInt}
\caption{Integer rectangle $a \times b$:
dome faces are  trapezoids and triangles.
}
\figlab{RectInt}
\end{figure}
%%%%%%%%%%%%%%%%%%%%%%%%%%%%%%%%%Figure End
%see Fig.~\figref{Rect_3x1_roof}.
%%%%%%%%%%%%%%%%%%%%%%%%%%%%%%%%%%Figure Begin
\begin{figure}[htbp]
\centering
\includegraphics[width=0.75\textwidth]{Figures/Rect_3x1_roof}
\caption{``Roof'' dome over a $3 \times 1$ rectangle.}
\figlab{Rect_3x1_roof}
\end{figure}
%%%%%%%%%%%%%%%%%%%%%%%%%%%%%%%%%Figure End

%Let $P_n$ be an equiangular convex polygon with integral edge lengths..
\subsection{Main Theorem}
%
\begin{theorem}
\thmlab{EquiAng}
\label{thm:main}
%\Anna{Equiangular implies convex, so I've removed convex}
(a)~The only equiangular %convex 
polygons
that can be domed have $n$ vertices, where $n \in \{3,4,5,6,8,10,12\}$;
for each such $n$, any regular integer $n$-gon  
can be domed.
(b)~Moreover, 
for $n=3,4,5,6$ every 
%convex 
equiangular integer polygon can be domed, and for $n=8,10,12$, an 
%convex 
equiangular integer $n$-gon can be domed iff the %odd edge lengths
lengths of the odd-numbered edges 
are equal and the even edge lengths are those of an equiangular $\frac{n}{2}$-gon.
\end{theorem}
%\begin{proof}
%\end{proof}


For small $n$, edge-length conditions for an equiangular integer polygon are known~\cite{ball2002equiangular}: 
%Anna: Ball gives conditions for n=5,6 (plus other results)
for $n=4$ these are rectangles; for $n=5$ there is only the regular pentagon; and for $n=6$ the edge lengths must be integers $a,b,c,a',b',c'$ with $a-a' = b'-b = c-c'$ (a 6-sided polyiamond);
see ahead to Fig.~\figref{Polyiamond6}.

Part~(a) of Theorem~\ref{thm:main} is proved in Sections~\ref{sec:RegPoly} and~\ref{sec:part-a}.
Part (b) is proved in Section~\ref{sec:main-part-b}.
We have established several results beyond the main theorem (for example
%, that all polyiamonds, equiangular or not, are domeable, and 
that there is no domeable polygon with $25$ or more vertices)---see Section~\ref{sec:further-results}.


%\hide{
%\noindent
%\Anna{This should be modified:}
%We just saw in Lemma~\lemref{Rect} that in fact all equiangular integer
%polygons of $n=4$ vertices---rectangles---can be domed. 
%\anna{This is not true for the other allowed values of $n$; we give the precise edge length conditions (for part (b) of the theorem) in Section~\ref{sec:further-results}. In that section we also mention some results beyond Theorem~\ref{thm:main}.}
%%However, it is not the case that every equiangular $n$-gon 
%%for $n$ in the claimed list can be domed.
%% Anna: the following is now incorporated into the Theorem:
%%But for each of those $n$, there is an equiangular $n$-gon that can be domed.
%} % end hide

\subsection{Glazyrin and Pak}
The source of our work derives from a paper by
Glazyrin and Pak: ``Domes over Curves''~\cite{glazyrin2022domes},
which answers a question posed by Richard Kenyon in 2005.\footnote{%
\url{https://gauss.math.yale.edu/~rwk25/openprobs/}.}
In~\cite{glazyrin2022domes}, a ``curve'' $P$ is a closed polygonal chain in $\R^3$,
and a dome is 
a PL-surface composed of unit equilateral triangles. 
They say that $P$ can be \emph{spanned}
if there is such a surface
whose boundary is $\partial P$.
We note the following two differences with our definitions:
\begin{enumerate}[(1)]
\item Our $P$ is a 2D convex polygon; theirs is a 3D possibly self-intersecting polygonal chain.
\item Our dome $\D$ is embedded (non-self-intersecting) and convex.
Their PL-surface is (in general) nonconvex, immersed, and self-intersecting.
\end{enumerate}
Under their conditions, they show that 
certain nonplanar rhombi cannot be spanned, which answers Kenyon's question in the negative.\footnote{%
Recent work~\cite{anan2024moduli}
extends the Glazyrin-Pak negative result to
show that ``generic'' integer polygons cannot be spanned. See~\cite{CobordismDomes} for a further generalization.
% \JOR{The first paper is a clear extension. The second paper generalizes to ``a polyhedron with multiple boundary components.''}
}
%JOR{``Our method is to generalize [1, Theorem 4.3] to allow for graph surfaces of arbitrary genus. We rephrase a cobordism as a sample polyhedron with multiple boundary components, and project down to smaller moduli space to account for any possible polygon to be chosen for the cobordism.''}
More interesting for our purposes, they prove that every planar regular polygon can be spanned
(their Theorem~1.4).
In contrast, our Theorem~\thmref{EquiAng} says that the regular $7$-, $9$-, and $11$-gons cannot
be domed, nor can any regular $n$-gon for $n>12$.
And here ``regular'' can be 
strengthened to ``equiangular.''
Compared to their results,
our conditions constrain the geometry and limit what can be domed.

% Anna: this was mentioned just after the Main Theorem
%Although we have established several results beyond Theorem~\thmref{EquiAng},
%here we focus on equiangular $P$.
%, only mentioning other results for future papers.

\section{Domed Regular Polygons}
%\label{sec:RegPoly}
\seclab{RegPoly}
We prove one part of Theorem~\ref{thm:main}(a)
by exhibiting domes over regular integer $n$-gons, for $n \in \{3,4,5,6,8,10,12\}$.
We will use $\bar{P}_n$ to denote a regular integer $n$-gon.
%(and later, $P_n$ for an equiangular $n$-gon).

\begin{itemize}
\item $3,4,5$ : $\bar{P}_n$ for $n=3,4,5$ can each be domed by a pyramid:
Fig.~\figref{Pyramids345}.
\item $6$ : Hexagonal antiprism: Fig.~\figref{SliceFigs}(a).
\item $8$ : A slice through a
gyroelongated
square diprism: Fig.~\figref{SliceFigs}(b).
\item $10$ : A slice through an icosahedron: Fig.~\figref{SliceFigs}(c).
\item $12$ : A slice through a hexagonal antiprism: Fig.~\figref{SliceFigs}(d).
\end{itemize}

A few remarks.
The pyramid pattern for $n=3,4,5$ cannot be extended to $\bar{P}_6$,
for that would result in a doubly-covered hexagon, not a
%delta 
dome by our definition.
For $n=8,10,12$, we show $\bar{P}_n$ as a slice of a convex polyhedron,
with the dome the upper half of the surface.
But we 
establish
that not every doming
of an equiangular polygon derives from a slice, see Section~\secref{NotSlice}.



%see Fig.~\figref{Pyramids345}.
%%%%%%%%%%%%%%%%%%%%%%%%%%%%%%%%%Figure Begin
\begin{figure}[htbp]
\centering
\includegraphics[width=1.0\textwidth]{Figures/Pyramids345}
\caption{Pyramids over $\bar{P}_n$, $n=3,4,5$.}
\figlab{Pyramids345}
\end{figure}
%%%%%%%%%%%%%%%%%%%%%%%%%%%%%%%%%Figure End


%see Fig.~\figref{SliceFigs}.
%%%%%%%%%%%%%%%%%%%%%%%%%%%%%%%%%Figure Begin
\begin{figure}[htbp]
\centering
\includegraphics[width=1.0\textwidth]{Figures/SliceFigs}
\caption{Examples of $\bar{P}_n$ domes for $n=6,8,10,12$.
%\Anna{Could the label font size (``(a) $n=6$'') be reduced ...} %to match the text font size in the caption (``$n=6,8,10,12$'')?}}
%\JOR{I'd rather not. Caption font is so small!
%But I reduced it a bit.}
}
\figlab{SliceFigs}
\end{figure}
%%%%%%%%%%%%%%%%%%%%%%%%%%%%%%%%%Figure End

%see Fig.~\figref{IcosaSlicesDifferent}.
%%%%%%%%%%%%%%%%%%%%%%%%%%%%%%%%%Figure Begin
\begin{figure}[htbp]
\centering
\includegraphics[width=0.85\textwidth]{Figures/IcosaSlicesDifferent}
%\includegraphics[width=1.0\textwidth]{Figures/IcosaSlicesDifferent}
\caption{(a)~A different dome over $\bar{P}_5$.
(b)~Doming an equiangular decagon with edge lengths alternating $1,3$.
}
\figlab{IcosaSlicesDifferent}
\end{figure}
%%%%%%%%%%%%%%%%%%%%%%%%%%%%%%%%%Figure End

Figures~\figref{Pyramids345} and~\figref{SliceFigs}
%These figures 
show one way to dome each regular polygon $\bar{P}_n$, but there are other solutions.
%These figures show that each of those regular polygons $\bar{P}_n$ can be domed, but says nothing
%about the number of different ways to dome.
%\Anna{Do we say more about the different ways to dome?}
For example, $\bar{P}_5$ can be domed by a low slice through the icosahedron
as shown in
Fig.~\figref{IcosaSlicesDifferent}(a).
And again, these figures illustrate regular polygons, special cases of equiangular polygons.
To give a hint of the further possibilities, Fig.~\figref{IcosaSlicesDifferent}(b) shows 
how to dome an equiangular decagon
$P_{10}$ whose edge lengths alternate $1$ and $3$.



%\JOR{Agreed on Notation: superfluous now.}
% \subsection{Notation}
% \seclab{Notation}
% % Anna: we don't really need the following sentence.
% %Here we fix notation for the proofs to follow; consult as needed.
% %The reader may skip these and return when needed.
% %\JOR{It is so convenient to use $v_i$ for a dome vertex and $b_j$ for a base vertex,
% %that I would like to redefine our coauthor notation as below.}

% %\Anna{Let's double-check which of these we actually use.}
% %\JOR{All used.} \Anna{$P_n$??}
% \Anna{I really think this notation section has become redundant.  $\bar{P}_n$ is only used in Section 2 and is defined there.  The others are used only in Sections 3 and 4 and are defined there. $P_n$ is not used at all.}
% \begin{itemize}
% \item $P$: Base polygon. $P_n$: Equiangular polygon.  $\bar{P}_n$: Regular polygon.
% %\item 
% $b_i$ : base vertex.
% \item $\D$: Dome.
% %\item 
% % $v_j$ : dome vertex.
% % \Anna{I don't think we use $v_j$}
% \item $V_k$ : The number of dome vertices 
% %of degree-$k$; so 
% % Anna: it's dangerous to call this degree $k$
% with $k$ incident equilateral triangles.
% %\item $B_k$ : A base vertex type with $k$ incident equilateral dome triangles.
% \item $\b$ : the base angle in an equiangular polygon.
% %\item 
% $\b_i$ : the base angle at vertex $b_i$ (when not all equiangular). 
% % \Anna{Check if we really need $\b_i$ and $b_i$.}
% % \JOR{It is used later.}
% \end{itemize}

%\clearpage
\section{Proof of Theorem~\thmref{EquiAng}(a): Restrictions on $n$}
\label{sec:part-a}
In this section we complete the proof of
the first half of Theorem~\thmref{EquiAng}:
The only equiangular %convex 
$n$-gons 
that can be domed have $n \in \{3,4,5,6,8,10,12\}$.
%\Anna{I deleted the blanket statement.}
%Throughout this section $\cal D$ is a dome over a base equiangular $n$-gon $P_n$, $n \ge 6$. 
% \JOR{I know this enables shortening the lemma
% statements, but I would prefer that each lemma
% state a truth independent of such a blanket
% statement.}
For a dome over an equiangular $n$-gon, $n \ge 6$, we use the following steps: 
%
\begin{enumerate}[(1)]
\item Each base vertex has three incident dome triangles.

\item Curvature constraints imply that 
the number of (non-base) dome vertices is at most 6.

\item Of the $n$ dome faces incident to base edges, at least half tilt toward the outside of the base
and no two of these faces share a dome vertex.
%and have  
%\anna{at least one ``overhanging'' dome vertex
%uniquely associated with the base edge.}
% \Anna{Trying to be more precise than ``private''. Our previous wording of ``unique'' could be interpreted as saying that each base edge has a single private vertex.}
%a ``private'' dome vertex. 
Furthermore, for $n$ odd we strengthen this claim to \emph{all} dome faces incident to base edges. 

\item Thus, since there are at most 6 dome vertices, 
%there are at most 12 base edges so 
$n \le 12$, and for $n$ odd, 
there are no solutions for $n \ge 6$.
%Anna: changing this to match the re-worded proofs
%$n \le 5$.
\end{enumerate}

Note that the base angle $\b$ of   
an equiangular $n$-gon
%is $\frac{n-2}{n} \pi$ Jayson:
is $\frac{n-2}{n} 180^\circ$
so if $n \ge 6$, then every base angle is $\ge 120^\circ$. 
%is $\ge 120^\circ$. 
This weaker assumption 
on a domeable convex 
$n$-gon
is enough for most of our argument.


\begin{lemma}
\lemlab{3Triangles}
\label{lem:3Triangles}
If a 
%\anna{[convex]}
base vertex  $b_i$ has base angle 
%\anna{$120^\circ \le \b_i < 360^\circ$}
$\b_i \ge 120^\circ$,
then it is incident to three dome triangles. %vertices.
%In particular, }
%in a dome over an equiangular $n$-gon $P_n$, $n \ge 6$,
%each base vertex 
% Anna: I removed $b_i$ from the lemma statement since it's not used in the proof
%$b_i$ 
%is incident to three dome triangles.
\end{lemma}
\begin{proof}
%The base angle $\b = \frac{n-2}{n} \pi$ of $P_n$
%is $\ge 120^\circ$.
%Therefore a base vertex  
Base vertex $b_i$ cannot 
be incident to just one
or two triangles, otherwise the total face angle is $\le 120^\circ$, which
is not enough to span $\b_i$.
%A base vertex 
Vertex $b_i$ cannot
be incident to four triangles,
because 
%$\b_i + \frac{4}{3} \pi \ge 2 \pi$, Jayson:
$\b_i + 240^\circ \ge 360^\circ$,
and similarly for 
five (or more) triangles.
%
%\hide{
%The base angle $\b = \frac{n-2}{n} \pi$ of $P_n$
%is strictly greater than $120^\circ$ for all $n \ge 7$.
%%(For $n=7$, $\b = \frac{5}{7} \pi \approx 129^\circ$.)
%%128.571
%Therefore a base vertex $b_i$ could not 
%be incident to just one
%triangle, or two 
%triangles, 
%because the total face angle is $\le 120^\circ$, which
%is not enough to span $\b$.
%Vertex $b_i$ could not 
%be incident to four triangles,
%because $\b + \frac{4}{3} \pi \approx 369^\circ 
%> 2 \pi$, and similarly for 
%five triangles.
%%368.571
%%
%%So each $b_i$ of $P_n$ must be incident to exactly three dome triangles.
%} % end hide
\end{proof}

%\noindent
%An immediate consequence is that the base curvature is $2\pi$:
%\noindent
From this we can analyze the base curvature:

\begin{lemma} 
\lemlab{BaseCurv}
If every base vertex $b_i$ is incident to three dome triangles,
then the sum of the curvatures at the base vertices is $2\pi$.
% \JOR{I put back the 3-triangle condition, which is
% I believe necessary.}
% \Anna: Right, since we've removed the blanket conditions
\end{lemma}
\begin{proof}
Let $\b_i$ be the angle of $P$ at vertex $b_i$.
Then the curvature at $b_i$ is $\o_i = 2 \pi - (\b_i + \pi)$, where the final $\pi$ term follows from
the assumption that $b_i$ is incident to three triangles.
Recalling that $\sum_i \b_i = \pi(n-2)$ for any simple polygon,
we have:
%\Anna{one line + fixed typos}
%\anna{
\begin{align*}
\sum_i \o_i  \ = \ \sum_i (\pi - \b_i) 
\ = \  n \pi - \sum_i \b_i 
 \ = \  n \pi - \pi (n-2) 
 \ = \  2 \pi.
\end{align*}
\end{proof}

%\hide{
%\begin{align*}
%\sum_i \o_i & = \sum_i (\pi - \b_i) \\
%%\sum_i \o_i & = \sum_i \pi - \o_i \\
%& = n \pi - \sum_i \b_i \\
%%& = n \pi - \sum_i \o_i \\
%& = n \pi - \pi (n-2) \\
%& = 2 \pi
%\end{align*}}



%\bigskip
%\noindent
%Note that the the sum could be greater or less than $2\pi$ without
%the three-triangles assumption.

% \hide{
% \Anna{We can eliminate this for space saving:}
% \begin{corollary}
% \corlab{GaussBonnet}
% Under the three-triangles assumption (in Lemma~\lemref{BaseCurv}),
% the total curvature of the dome vertices is $2 \pi$.
% \end{corollary}
% %\begin{proof}
% %This follows immediately from the Gauss-Bonnet theorem, which
% %says that the total curvature is $4 \pi$.
% %\end{proof}
% } % end hide
%

%\Anna{I'm tempted to move the following argument inside the next lemma so we can be clear about the hypotheses. But the trouble is our assumption about no $V_3$ vertices.  Hmmm.  I'll think about this.}



\begin{lemma}
\lemlab{6dome}
\label{lem:6dome}
If $\D$ is a dome over a convex polygon $P$ that has all angles $\ge 120^\circ$, then $\D$ has at most $6$ dome vertices.
% \Anna{we can word this more simply since we no longer have to cut off $V_3$ vertices.}
\end{lemma}
\begin{proof}
Let 
$V_3, V_4, V_5$ be the number of (non-base) dome vertices with $3,4,5$ incident triangles,
respectively.
% \JOR{Reviewer~4 thinks $V_k$ are not clearly defined. Seems clear to me.}
%If $\D$ has a $V_3$ vertex $v$, then $v$ is the apex of a rigid regular tetrahedron
%which we can remove without changing $P$.  
%Thus, $P$ has a dome $\D'$ with no $V_3$ vertex.
%We prove that $\D$ has at most 6 dome vertices.
% \JOR{Reviewer~1 argues that we need not set $V_3=0$. Seems correct.}
%
By the Gauss-Bonnet theorem, the total curvature of a convex polyhedron is $4\pi$. 
The curvature of a $V_k$ vertex is $2\pi - k \frac{\pi}{3}$.
By Lemmas~\ref{lem:3Triangles} and~\lemref{BaseCurv} 
(this is where we use the assumption that all base angles are $\ge 120^\circ$)
the curvature at the base vertices of any dome over $P$ is $2\pi$. 
%Since a $V_i$ vertex has curvature $2\pi - i\frac{\pi}{3}$ the curvature of the dome vertices is
Thus, in units of~$\pi$:
\begin{align*}
V_3 +  \frac{2}{3} V_4 + \frac{1}{3} V_5 &= 2 \; .
%3 V_3 +  2 V_4 +  V_5 &= 6 \;. 
\end{align*}
%Since $\D'$ has no $V_3$ vertices, this becomes:
%We may assume that $V_3=0$:
%Any $V_3$ vertex is the apex of a rigid regular tetrahedron
%which we can remove without %changing $P_n$.
%The equation above becomes:
%\begin{align*}
%2 V_4 + V_5 = 6 \; .    
%\end{align*}
%Therefore
Thus
% Anna: changed just to put the equation on one line.  If we use fullpage, this change isn't needed.
the number of dome vertices of $\D$ is 
$V_3 + V_4 + V_5 \le 3 V_3 +  2 V_4 +  V_5  = 6$. 
%
\end{proof}

% \Anna{Note the weirdness that the qed box is on the left.}
%\JOR{Left it as-is.}

% \hide{
% \subsection{Dome Vertices Curvature}


% \seclab{DomeCurv}
% %\bigskip
% \noindent
% We will now use this fact to constrain the possible number of dome vertices
% of each possible degree.
% Let $V_3, V_4, V_5$ be the number of dome vertices with $3,4,5$ incident triangles.
% %(Note: A $V_k$ vertex might not have $k$ incident nonflat edges; some of the $k$ edges could be flat.)


% First we give a high-level view of the next three steps in the argument,
% followed by detailed justifications.


% %Each $V_3, V_4, V_5$ vertex has curvature $\pi, \frac{2}{3} \pi, \frac{1}{3} \pi$ respectively.
% %Therefore, to satisfy Corollary~\corref{GaussBonnet}, we must have
% %\begin{align}
% %V_3 +  \frac{2}{3} V_4 + \frac{1}{3} V_5 &= 2 \\
% %3 V_3 +  2 V_4 + V_5 &= 6 \eqnlab{curv}
% %\end{align}


% \begin{enumerate}[(a)]
% %
% \item We may assume that $V_3=0$.
% Let $v$ be such a degree-$3$ dome vertex.
% Then $v$ is the apex of a rigid regular tetrahedron.
% We can remove that tetrahedron---a kind of ``anti-augmentation''---and glue it back
% if the remainder is domed.
% %
% \item 
% The Gauss-Bonnet theorem requires, in units of $\pi$:
% \begin{align}
% V_3 +  \frac{2}{3} V_4 + \frac{1}{3} V_5 &= 2 \\
% 2 V_4 + V_5 &= 6 \eqnlab{curv}
% \end{align}
% where we have set $V_3=0$.
% This is our first equation in the $V_k$ variables.
% %
% %\item The dome vertices cannot consume all of the $4 \pi$ curvature---some must be
% %left for the base vertex curvatures, each of which is non-zero: 
% %\begin{align}
% %\frac{2}{3} V_4 + \frac{1}{3} V_5 & < 4 \\
% %2 V_4 + V_5 & < 12  \eqnlab{curv}
% %\end{align}
% %So, in particular, a $V_4,V_5$ dome can contain at most $11$ vertices.
% %
% \item If $|P|=n$, then
% $$2 V_4 + 3 V_5 + 2 \ge n \; .$$
% This follows from Lemma~\lemref{Connected} below,
% and again setting $V_3=0$.
% %
% \end{enumerate}

% \Anna{The special case we need is incorporated into Lemma 3.1}
% \begin{lemma}
% \lemlab{NonFlat2Base}
% For any base $P$ with each base vertex angle $\b_i > 120^\circ$,
% $b_i$ must have at least one incident nonflat edge.
% \end{lemma}
% \hide{
% \begin{proof}
% Let $b$ be a base vertex with base angle $\b < \pi$.
% Consider the number of dome triangles incident to $b$.
% %Recall by our definition, no dome triangle can lie in the plane of the base.
% %, and not lying
% %in the plane of the base $P$, which would violate either (a)~the domed base being a polyhedron,
% %or (b)~would extend the base beyond $B$.
% It cannot be that 
% either only $1$ or $2$ triangles incident to $b_i$, because of the assumption that $\b_i > 120^\circ$.
% If the number is $3$, then if both incident edges are flat,
% then $\b=\pi$ and $b$ is not a vertex. So one of the incident edges must be nonflat.
% If the number if $4$ or $5$, there must be $3$ or $4$ incident nonflat edges respectively.
% \end{proof}
% }%hide

% \Anna{We don't need this; see below}
% \begin{lemma}
% \lemlab{Connected}
% The graph $G$ of nonflat edges of the dome, excluding the base edges, 
% \Anna{and excluding base vertices -- isn't that how we're using it?} is connected.
% \end{lemma}
% %AL:
% %The same argument as in Balinski’s theorem should do it.  Orient the dome with the base horizontal and let $P$ be the horizontal supporting plane of the highest vertices of the dome.  The part of the dome on $P$ is connected since it’s a vertex, edge, or polygon.
% %To show that the graph of the dome minus the base is connected it suffices to show that there is a path from any dome vertex $v$ to a vertex on $P$  If $v$ has no edge going upward then $v$ already lies on $P$  Otherwise there is an upward edge from $v$ to a vertex $u$ and by induction there’s a path from $u$ to a vertex on $P$
% \begin{proof}
% Orient the dome with the base horizontal and let $\Pi$ be the horizontal supporting plane of the highest vertices of the dome.  The part of the dome on $\Pi$ is connected since it’s a vertex, edge, or polygon.

% To show that $G$ minus the base is connected, it suffices to show that there is a path from any dome vertex $v$ to a vertex on $\Pi$.  
% If $v$ has no nonflat edge heading upward then $v$ already lies on $\Pi$.

% %Suppose $v$ has only a flat edge $e_2$ heading upward.
% %Then $v$ must be of degree-$5$, because degree-$3$ and degree-$4$ vertices have only
% %nonflat incident edges. But if $e_2$ slopes upward,
% %its two adjacent edges $e_1$ and $e_3$ must also slope upward, because
% %the angle between $e_1$ and $e_3$ is $120^\circ$

% Otherwise there is an upward nonflat edge from $v$ to a vertex $u$, 
% and by induction there’s a path from $u$ to a vertex on $\Pi$.
% \end{proof}

% \noindent
% An alternative viewpoint is that the $1$-skeleton of a convex polyhedron, 
% minus the edges of one face, is connected.

% \medskip

% \Anna{We don't need this either.}
% \noindent
% We may now obtain a second equation with variables $V_k$:

% \begin{lemma}
% \lemlab{nBoundEq}
% If the base $P$ is an $n$-gon, 
% \Anna{Stronger assumption needed in order to use Lemma 3.4 in the proof.}
% then
% \begin{align}
% V_3 + 2 V_4 + 3 V_5 + 2 \ge n \; . \eqnlab{gen}
% \end{align}
% \end{lemma}
% \hide{
% \begin{proof}
% %\JOR{From Erik's post.}
% %https://coauthor.csail.mit.edu/compgeom/m/endaQJCfAkXkSC29o
% By Lemma~\lemref{Connected},
% the graph $G$ of nonflat edges of the dome is connected.
% $G$ has a spanning tree, which consumes $2(V-1)$ total degree
% among the  $V = V_3+V_4+V_5$ dome vertices.
% So the total number of nonflat edges is
% \begin{align}
% \leq & \; 3 V_3 + 4 V_4 + 5 V_5 - 2(V-1) \\
% = & \; V_3 + 2 V_4 + 3 V_5 + 2 \;. 
% \end{align}
% The number of base vertices $n$ is at most this quantity, because
% by Lemma~\lemref{NonFlat2Base}
% base vertex has at least one nonflat edge to the dome.
% \end{proof}
% }%hide


% %Namely, while we incorrectly argued that the nonflat edges of the dome are 2-connected, I think they must be connected: I believe multiple connected components in the dome would force nonconvexity between these multiple “peaks”.  So the nonflat edges in the dome have a spanning tree, which consumes $2(v-1)$ total degree among the $v = v_3+v_4+v_5$ dome vertices.  Thus if we sum the nonflat-edge degree of the dome vertices, we get
% %$$ \leq 3 v_3 + 4 v_4 + 5 v_5 - 2(v-1) = v_3 + 2 v_4 + 3 v_5 + 2, $$
% %and $n$ must be at most this because every base vertex has at least one nonflat edge to the dome


% %\subsection{Solutions to the Curvature Equations}
% There are only four 
% solutions of the two equations in nonnegative integers,
% Eqs.~\eqnref{curv} and~\eqnref{gen}, shown in the Appendix.
% These four cases lead to this conclusion:

% \hide{
% \begin{align}
% 2 V_4 + V_5 & = 6 \\
% 2 V_4 + 3 V_5 + 2 & \ge n  %\eqnlab{gen}
% \end{align}

% \noindent
% The solutions to Eqs.~\eqnref{curv} and~\eqnref{gen} are gathered in the following table.

% \noindent
% \begin{center}
% \begin{tabular}{| c || c | c |}
% \hline 
% $n$ & $V_4$ & $V_5$ \\
% \hline\hline
% $n \ge 21$ & $\varnothing$ & $\varnothing$ \\
% \hline
% $n \in [3,20]$ & $0$ & $6$ \\
% \hline
% $n \in [3,16]$ & $1$ & $4$ \\
% \hline
% $n \in [3,12]$ & $2$ & $2$ \\
% \hline
% $n \in [3,8]$ & $3$ & $0$ \\
% \hline\hline
% \end{tabular}
% \end{center}

% One immediate consequence is that no equiangular polygon of more than $20$ vertices can be domed.
% However, the equations reflect the equiangular assumption (in Lemma~\lemref{3Triangles}), so
% they do not establish a bound for all base polygons.
% Nor do the equations imply there are realizations for each solution;
% rather they represent constraints on any realization. \JOR{Correct?}

% %A few examples: %$V_4=0$ and $V_5=6$ 
% %\begin{itemize}
% %\item 
% %$V_4,V_5 = 0,6$ is realized
% %by the pentagon in Fig.~\figref{PentagonIcosaDome},
% %\item 
% %by the hexagon in  Fig.~\figref{HexagonalAntiprism},
% %\item
% %by the decagon in Fig.~\figref{DecagonSlice},
% %\item
% %and by the dodecagon in Fig.~\figref{DodecagonSlice}.
% %\item
% %%$V_1=1$ and $V_5=4$ 
% %$V_4,V_5 = 1,4$ is realized by the octagon in Fig.~\figref{OctagonSlice}.
% %\item
% %$V_4,V_5 = 2,2$ can be realized with a square base and a two-triangle top.
% %%Polydron image in coauthor.
% %%\JOR{Any more? That's probably more than enough.}
% %\end{itemize} 


% Although we can draw conclusions by analyzing the cases in this table individually,
% it turns out that for our immediate goal we only need this conclusion:
% }%hide
% %Although we can draw conclusions by analyzing the four cases individually,
% %it turns out that for our immediate goal we only need this conclusion:
% %These four cases lead to this conclusion:
% \begin{lemma}
% \lemlab{6verts}
% For an equiangular base $P_n$, $n \ge 7$, there can be at most $6$ dome vertices.
% \end{lemma}
% \begin{proof}
% \Anna{We had a slick proof of this: since we assume no $V_3$ vertices, the total number of dome vertices is $V_4 + V_5 \le 2V_4 + V_5$ and $2V_4 + V_5=6$ by equation (2).  We don't need the spanning tree argument at all.}

% The equiangular assumption and and $n \ge 7$ lead, via Lemma~\lemref{3Triangles} to
% three triangles per base vertex.
% This then leads via Lemma~\lemref{BaseCurv} to Corollary~\corref{GaussBonnet}:
% the dome curvature sum is $2\pi$.
% This leads to Eq.~\eqnref{curv}.
% The two nonflat edge lemmas (Lemmas~\lemref{NonFlat2Base} and~\lemref{Connected})
% lead to Lemma~\lemref{nBoundEq} and Eq.~\eqnref{gen}.
% Solving the two equations leads to %the displayed table, and yields 
% the claim of the lemma.
% \end{proof}
% } %end hide

% \hide{
% %\clearpage
% \subsection{Normals and Projection}
% This completes the first four steps of the proof.
% The last two steps introduce further geometric constraints that ultimately
% show the impossibility of odd $n \ge 7$, and then even $n \ge 14$.
% Very roughly, the ``$\pm$Normals Lemma''~\lemref{Normals} shows that if a base vertex
% has exactly three incident dome triangles (which we know is the case
% from Lemma~\lemref{3Triangles}),
% then the faces incident to a base edge tilt toward the outside of the base
% \Anna{that's too strong a statement unless we add that $n$ is odd}.
% This then leads by a projection argument and Lemma~\lemref{6verts} to concluding 
% that $n$ cannot be large.
% } % end hide

%\subsubsection{$\pm$Normals Lemma}
%\seclab{Normals}
%\JOR{Based on ``Proof of missing claim."
%\url{https://coauthor.csail.mit.edu/compgeom/m/ix8tk5zqWNtE7GaCX}}


%\JOR{Removed notation list experimentally to see if it can be done in-line.}
%\subsubsection{Notation-II}
%\seclab{Notation-II}
%
%\begin{itemize}
%%\item $P$: the base polygon, in the $xy$-plane.
%%\item $b_i$ : base vertex.
%\item $t_1, t_2, t_3$: the three triangles incident to base vertex $b_2$.
%\item $a_1$: apex of $t_1$.
%\item $a_2$: apex of $t_3$.
%\item $t_1 = b_1 b_2 a_1$,
%$t_2 = b_1 a_1 a_2$,
%$t_3 = b_2 b_3 a_2$.
%\item $n_1,n_2,n_3$: normal vectors to $t_1, t_2, t_3$.
%\item $n_0 = (0,0,-1)$, the downward normal to the base.
%\item $\b$: Angle of the base at $b_2$.
%\end{itemize}
%See Fig.~\figref{Notation_t1t2t3}.
%%%%%%%%%%%%%%%%%%%%%%%%%%%%%%%%%%Figure Begin
%\begin{figure}[htbp]
%\centering
%\includegraphics[width=0.5\textwidth]{Figures/Notation_t1t2t3}
%\caption{Notation.}
%\figlab{Notation_t1t2t3}
%\end{figure}
%%%%%%%%%%%%%%%%%%%%%%%%%%%%%%%%%%Figure End


\subsection{Face Normals and Private Dome Vertices}
\seclab{NormalsPrivate}

% \Anna{I reworked this section.  In the EuroCG version, the main lemma (Lemma~\ref{lem:even-n}) was proved with a sequence of nested lemmas, so the proof was the whole rest of the section.  In the first arXiv version, an ``end\{proof\}'' was added too early (you can see this in the old grayed out version below).  Anyhow, now that we give complete proofs, the nested lemmas are too confusing.  }

In this section we complete the proof of Theorem~\thmref{EquiAng}(a) by proving:

\begin{lemma}
\label{lem:12-and-6}
If $P$ is a domeable convex $n$-gon with all angles $\ge 120^\circ$, then $n \le 12$. 
Furthermore, there is no domeable equiangular $n$-gon for odd $n$ $ \ge 6$.   
\end{lemma}

To prove this, we 
consider a dome over $P$ and
follow step (3) which states that, of the $n$ dome faces incident to base edges, at least half of them [and for odd $n$, all of them] tilt toward the outside of the base
and no two of these faces share a dome vertex.
%each have a ``private'' dome vertex.  
We begin by making this more formal.
Orient the dome with the base in the horizontal $xy$-plane. A dome triangle/face has an \defn{upward normal} if its outer normal has a 
positive $z$-component, and a \defn{downward normal} if its normal has a negative $z$-component.
%(This formalizes ``tilting towards the outside'').
Let $d$ be the number of dome faces with downward normals that are incident to base edges. 
Each of these faces %has 
is incident to
at least one dome vertex, and we prove that no two share a dome vertex.
%uniquely associated with the base edge,
%can be associated with a unique dome vertex, 
This implies that $d \le 6$ (since there are at most $6$ dome vertices).  We then prove that  
$d \ge n/2$, which implies that $n \le 2d \le 12$, and we prove that for odd $n$, 
$d \ge n$, which implies that $n \le d \le 6$.

% \old{It remains to show step (3): that, of the $n$ dome faces incident to base edges, at least half of them [and for odd $n$, all of them] tilt toward the outside of the base and each have a ``private'' dome vertex.
% % \JOR{Reviewer~4 thinks this is not clearly defined. Here's Pepa's revision:}
% % \pepa{Here we say that a dome vertex $v$ is \defn{private} to a base edge $e$ if $e$ is the unique base edge whose non-base face is incident to $v$.}
% %\Anna{How about:} 
% %
% We say that a dome vertex $v$ is \defn{private} if 
% there is a unique dome face incident to $v$ and to a base edge.
% %\Anna{What do we mean exactly?  That there is exactly one dome face incident to $v$ and to a base edge, right?  It might be clearer to say that.}
% Orient the dome with the base in the horizontal $xy$-plane. A dome triangle/face has an \defn{upward normal} if its normal has a 
% positive $z$-component, and a \defn{downward normal} if its normal has a negative $z$-component.
% (This formalizes ``tilting towards the outside'').
% }

We first analyze upward and downward normals of the faces at a base vertex.

\begin{lemma}[$\pm$Normals]
\lemlab{Normals}
\label{lem:Normals}
% \Anna{Whoops, the proof uses the stronger assumption that the base vertex has angle $\ge 120^\circ$, so we need to add it as a hypothesis here.}
Consider a base vertex $b_i$ with base angle $\ge 120^\circ$.  Suppose the three dome triangles incident to $b_i$
are $t_1, t_2, t_3$ where $t_1$ and $t_3$
%lie in dome faces $f_1, f_2, f_3$, where $f_1$ and $f_3$
are incident to the base edges at $b_i$
(possibly $t_2$ is coplanar with $t_1$ or with $t_3$, but not both).
%(possibly $f_1=f_2$ or $f_2=f_3$).  
See Fig.~\figref{3Tri_Overhead2D}.
Then 
$t_2$ has an upward normal 
and at least one of $t_1, t_3$
has a downward normal.
% \Anna{The notation $f_i$ is not consistent with the Appendix.  Change?}
% \JOR{We decided that's OK.}
\end{lemma}
%See Fig.~\figref{3Tri_Overhead2D}.
%%%%%%%%%%%%%%%%%%%%%%%%%%%%%%%%%Figure Begin
\begin{figure}[htbp]
\centering
%\includegraphics[width=0.75\textwidth]{Figures/3Tri_Overhead2D}
\includegraphics[width=0.6\textwidth]{Figures/3Tri_Overhead2D}
\caption{
Overhead view of three  triangles
incident to base vertex $b_i$.
Triangles with upward normals pink, downward normals blue. (a)~Both $t_1$ and $t_3$ downward.
(b)~Only 
$t_1$ downward.
}
\figlab{3Tri_Overhead2D}
\end{figure}
%%%%%%%%%%%%%%%%%%%%%%%%%%%%%%%%%Figure End

%The argument depends on the Gauss map, which we review before starting the proof.

%see Fig.~\figref{GSph_generic}.
%%%%%%%%%%%%%%%%%%%%%%%%%%%%%%%%%Figure Begin
\begin{figure}[htbp]
\centering
\includegraphics[width=0.8\textwidth]{Figures/GSph_generic}
%\includegraphics[width=1.0\textwidth]{Figures/GSph_generic}
\caption{The Gauss map for the vertex $v$.}
\figlab{GSph_generic}
\end{figure}
%%%%%%%%%%%%%%%%%%%%%%%%%%%%%%%%%Figure End


%\subsection{Gauss Map}
%\section{Gauss Map}
The argument depends on the Gauss map, which we first review.
We start with the Gauss map
% Anna: deferring $\S$ to clarify that what we start with is a generic $v$
%on the unit sphere $\S$ 
for the neighborhood of 
a generic vertex $v$ of
degree $4$,
as described for example in~\cite{biedl2007cauchy}.
%\Anna{$n_i$ bothers me.  ...} 
% But I see that it is one convention for face normals.  I also see it bolded, and I see $u_i$.  OK, I won't fuss about this.}
%\JOR{Let's leave it. Bold in the figs does not look great.}
The four face normal vectors $n_i$, $i=1,2,3,4$, map to points
on the unit sphere $\S$.
This mapping can be viewed as placing the unit normals $n_i$ at the origin $o$
of $\S$, so the tips of $n_i$ are points on $\S$.
See Fig.~\figref{GSph_generic}.
Vertex $v$ and the edges incident to $v$ each map to their set of normals.  Thus the edge $e_i$ between the faces with normals $n_i$ and $n_{i+1}$ maps to an arc on $\S$ from $n_i$ to $n_{i+1}$, and $v$ itself maps to a spherical quadrilateral on $\S$. If $e_i$ has dihedral angle $\d_i$, then the corresponding arc has length $\pi - \d_i$.
%The dihedral angles along the four edges incident to $v$ become
%arcs on $\S$, reflecting the rotation of one normal to the next across the edge they share.
%Dihedral angle $\d$ corresponds to an arc of length $\pi-\d$.
For example, a $\d_i$ near $\pi$ indicates two nearly coplanar faces and
so two nearly parallel normals, leading to a near zero arc length.
Each face angle $\a_i$ incident to $v$ is an angle between two edges,
which on $\S$ becomes an angle between the corresponding %dihedral 
arcs.
That angle is $\pi-\a_i$;
see, e.g.,~\cite{s-gpcpl-68}.
%\JOR{Reviewer~4 objects that Stoker didn't originate that, it was known long before.
%Likely correct, so now just ``e.g.''}
%That angle is $\pi-\a_i$,
%as was proved by Stoker~\cite{s-gpcpl-68}.
%\Anna{I suggest not repeating Stoker's argument ...}
%\JOR{Agreed}
%(so remove the next sentence and paragraph).}
% That the angle on $\S$ is the supplement of the face
% angle can be seen as follows, following Stoker~\cite{s-gpcpl-68}.

% Let the point $n_1$ on $\S$ be the normal to the face whose face angle
% at $v$ is $\a$.
% Suppose the two edges determining the face angle $\a$ at a vertex are $e_1, e_2$,
% at which the dihedrals are $\d_1, \d_2$.
% Those dihedrals map to arcs on $\S$ incident to $n_1$.
% Each arc is part of a great circle determined by a plane through the center $o$ of $\S$.
% If we translate the planes to pass through $v$, they are each orthogonal to edges $e_1, e_2$.
% If we imagine $\a$ shrinking, drawing $e_1, e_2$ together, then those planes
% approach parallelism, i.e., the angle between them on $\S$ is $\pi-\a$,
% with the arcs approaching collinearity.
% See Fig.~\figref{Supplement}.
% %%%%%%%%%%%%%%%%%%%%%%%%%%%%%%%%%Figure Begin
% \begin{figure}[htbp]
% \centering
% \includegraphics[width=0.4\textwidth]{Figures/Supplement}
% \caption{Face angle $\a$ incident to $v$ maps to angle $\pi-\a$ on $\S$.
% Blue lines are the great circle planes seen edge-on.}
% \figlab{Supplement}
% \end{figure}
% %%%%%%%%%%%%%%%%%%%%%%%%%%%%%%%%%Figure End

%\begin{lemma}[$\pm$Normals]
%\lemlab{Normals}
%If base vertex $b_2$ is incident to three triangles $t_1, t_2, t_3$, 
%and if the base angle $\b \ge 120^\circ$, then
%$n_2$ is an upward normal, and at least one of $n_1,n_3$ is a downward normal.
%\end{lemma}
\medskip

% \Anna{We really need Theorem-restate.  I worry we'll change one version of a lemma and forget to change its repetition.}
%\noindent
%Recall Lemma~\ref{lem:Normals}:
%Consider a base vertex $b_i$ with base angle $\ge 120^\circ$.  Suppose the three dome triangles incident to $b_i$
%are $t_1, t_2, t_3$ where $t_1$ and $t_3$
%are incident to the base edges at $b_i$
%(possibly $t_2$ is coplanar with $t_1$ or with $t_3$, but not both).
%Then 
%$t_2$ 
%has an upward normal 
%and at least one of 
%$t_1, t_3$
%has a downward normal.
%

\begin{proof} [Proof of Lemma~\lemref{Normals}]
Let $n_i$ be the normal of $t_i$, $i=1,2,3$ and let $n_0$ be the normal of the base.
The proof analyzes the Gauss map 
for vertex $b_i$;
see Fig.~\figref{GSph_Pent}.
Because each triangle $t_1, t_2, t_3$ has angle $60^\circ$ incident to $b_2$,
on $\S$ the angles at $n_1,n_2,n_3$ are each $\pi-60^\circ=120^\circ$.
The normal 
$n_0$ is at the south pole of $\S$.
Because 
the base angle at $b_i$ is $\ge 120^\circ$,
we know the angle at $n_0$ on $\S$ is $\le 60^\circ$.
We first prove that $n_2$ is an upward normal, i.e., it lies above the equator
of $\S$.

Construct on $\S$ the arc $x$ from $n_0$ to $n_2$ that lies inside 
the spherical
quadrilateral $n_0 n_1 n_2 n_3$. 
Note that  $|x| \le 180^\circ$  (where $|x|$ is the length of $x$) because the spherical quadrilateral is contained in a lune from south to north pole with angle $\le  60^\circ$ (the angle at $n_0$).
Thus, proving that $n_2$ lies above the equator is equivalent to proving that $|x| > 90^\circ$.


% Note this issue about |x| < 180 degrees:
%The issue is that $x$ being \emph{inside} the quadrilateral might force it to be the longer of the two arcs from $n_0$ to $n_2$.  But in that case (IF it could arise, which is can't) it's NOT enough to show $|x| > 90^\circ$ ($x$ might be very close to the south pole).  


%This fig. from coauthor helps: https://coauthor.csail.mit.edu/file/ix8tk5zqWNtE7GaCX

%the base angle is >= 120 so the angle at n_0 is <= 60.

%Maybe you could enhance Fig. 9 to show the sides of the lune as in Fig. 2 in the figure I pointed to? 



% This was Erik's comment in coauthor, now resolved above
%Note that $|x|$ could be $> 180^\circ$, i.e., $360^\circ$ minus its dihedral angle,
%but still it's enough to show $|x| > 90^\circ$.
%
Arc $x$ splits the $120^\circ$ angle at $n_2$ into two parts, $A,B$, with $A$ on the $n_3$ side. 
% \Anna{In order to have $A$ on the $n_3$ side, please switch $n_1$ and $n_3$ in the figure (the original in coauthor is correct).  In Section 4 we used the other ordering of $t_1, t_2, t_3$ (this was remarked on in coauthor), but maybe let's not make things consistent at this late stage.}
Without loss of generality, assume $A \ge B$; so $A \ge 60^\circ$.  We also split the angle at $n_0$ into two parts $C,D$, with $C$ on the $n_3$ side.  We have $0 < C,D \le 60^\circ$. 

Apply the basic cosine law for spherical triangles
to the triangle $n_3 n_2 n_0$:
%, see https://en.wikipedia.org/wiki/Spherical_trigonometry#Supplemental_cosine_rules

$$\cos 120^\circ = - \cos A \cos C + \sin A \sin C \cos |x| \;.$$

Rearranging and using $\cos 120^\circ = -\frac{1}{2}$ leads to

$$\cos |x| = \frac{- \frac{1}{2} + \cos A \cos C}{\sin A \sin C} \;.$$

Now $A \ge 60^\circ$ so $\cos A \le \frac{1}{2}$.  Also, $\cos C < 1$.  Thus the numerator is negative.  Because the %$\sin$
sine terms are positive, the denominator is positive.  Thus $\cos |x| < 0$, 
and so $|x| > 90^\circ$. This establishes the first claim of the lemma: $n_2$ is an upward normal.

Now we turn to showing at least one of $n_1,n_3$ lies below the equator.
Let $y$ be the arc $n_0 n_1$.
Again apply the cosine law:

$$\cos B = - \cos 120^\circ \cos D + \sin 120^\circ \sin D \cos |y| \;.$$

$$\cos |y| = \frac{\cos B - \frac{1}{2} \cos D}{\sin 120^\circ \sin D} \;.$$

Now $B \le 60^\circ$ so $\cos B \ge \frac{1}{2}$.  Also, $\cos D < 1$.  Thus the numerator is positive.  Because the %$\sin$ 
\anna{sine} 
terms are positive, the denominator is positive.  Thus $\cos |y| > 0$, so $|y| < 90^\circ$, i.e., $n_1$ is a downward normal. 
\end{proof}

\medskip
\noindent
The asymmetry derives from the assumption that $A \ge B$. So we draw no conclusion about $n_3$'s location on $\S$. And indeed it is possible for $n_3$ to be upward or downward.

%See Fig.~\figref{GSph_Pent}.
%%%%%%%%%%%%%%%%%%%%%%%%%%%%%%%%%Figure Begin
\begin{figure}[htbp]
\centering
\includegraphics[width=0.50\textwidth]{Figures/GSph_Pent}
%\includegraphics[width=0.75\textwidth]{Figures/GSph_Pent}
\caption{The Gauss map for base vertex $b_2$. Here $n_2$ is an upward normal and
both $n_1,n_3$ are downward.
% \Anna{Note clockwise $n_i$ here, but counterclockwise in previous figure.}
%\Anna{Please swap $n_1$ and $n_3$.}
}
\figlab{GSph_Pent}
\end{figure}
%%%%%%%%%%%%%%%%%%%%%%%%%%%%%%%%%Figure End



Since every vertex of $P$ must be incident to at least one dome face with a downward normal, we immediately obtain the following result.
Recall that 
$d$ is the number of dome faces with downward normals that are incident to base edges.


\begin{corollary}
If $P$ is a domeable convex $n$-gon with all angles $\ge 120^\circ$, then $d \ge n/2$.    
\end{corollary}


Going forward, we need the following implication of the above proof.


\begin{obs} 
\label{obs:down-face}
In Lemma~\ref{lem:Normals}, if 
%$f_3$ 
$t_1$ has a downward normal, then it cannot be coplanar with 
$t_2$
%$f_2$ 
(which has an upward normal), and thus the 
dome face containing 
$t_1$ has a face angle of $60^\circ$ at the base vertex $b_i$.
\end{obs}


\begin{lemma}
\label{lem:d-le-6}
If $P$ is a convex $n$-gon with all angles $\ge 120^\circ$, and $\cal D$ is a dome over $P$, then $d \le 6$.
\end{lemma}
\begin{proof}
Consider a downward pointing face $f$ incident to base edge $e$. By Observation~\ref{obs:down-face}, $f$
has $60^\circ$ face angles at the endpoints of $e$, 
so the dome vertices of $f$, when projected to the $xy$-plane, lie in the equilateral $xy$-triangle on edge $e$.
See Fig.~\figref{3Tri_Overhead2D}
and Fig.~\figref{ProjectionPrivate}.
%
These triangles are disjoint, so 
no two downward pointing faces share a dome vertex.
Since $\D$ has
at most $6$ dome vertices (by Lemma~\ref{lem:6dome}) 
we have $d \le 6$.    
\end{proof}


% \old{
% \begin{lemma}
% \label{lem:even-n}
% If $P$ is a domeable convex $n$-gon with all angles $\ge 120^\circ$, then $n \le 12$. 
% Furthermore, there is no domeable equiangular $n$-gon for odd $n$ $ \ge 6$.
% \end{lemma}
% %
% \begin{proof}
% For the second statement in the lemma, note 
% that an equiangular $n$-gon, $n \ge 6$ has all angles $\ge 120^\circ$. 
% so
% consider
% an $n$-gon
% $P$ with all angles $\ge 120^\circ$,
% and suppose $P$ has a dome $\D$.
% By Lemma~\ref{lem:6dome},
% %$\D$ has at most 6 dome vertices.
% We will prove that $n \le 12$, and derive a contradiction for odd $n$ $\ge 6$.
% %
% % \JOR{Changed to $>$. I know you are trying to match the
% % hypothesis that $n \ge 6$ and $\ge 120$, but
% % seems strange to use $\ge 6$ for an odd number.}
% %\anna{Because all angles of $P$ are $\ge 120^\circ$, Lemma~\ref{lem:6dome} implies that}
% %\Anna{Careful.  We need $n\ge 6$ to claim angles $\ge 120^\circ$.  Nothing is incorrect, but we need to write this more carefully.}
% %By Lemma~\ref{lem:6dome}, 
% %$P$ has a dome $\D$ with at most 6 dome vertices.  
% %
% Consider the $n$ faces of $\cal D$ incident to a base edge. 
% Let $d$ be the number of those faces with downward normals.
% By Lemma~\ref{lem:Normals}, $d \ge \frac{n}{2}$. 
% Furthermore, a
% downward pointing face $f$ incident to base edge $e$ 
% has $60^\circ$ face angles at the endpoints of $e$, 
% so it must include a dome vertex whose projection
% to the $xy$-plane lies in the equilateral $xy$-triangle on edge $e$.  
% See Fig.~\figref{3Tri_Overhead2D}
% and Fig.~\figref{ProjectionPrivate}.
% %
% Thus face $f$ has a private dome vertex. 
% Such a dome vertex is unique to base edge $e$ so we call it a ``private'' dome vertex. 
% %\JOR{Reviewer~4 thinks this is not clearly defined...}
% Thus 
% there are at least $d$  %$\frac{n}{2}$ 
% private vertices.
% Since $\D$ has
% at most $6$ dome vertices (by Lemma~\ref{lem:6dome}) 
% we have $d \le 6$.
% Now Lemma~\ref{lem:Normals} implies that $d \ge \frac{n}{2}$.  
% \end{proof}

% Now Lemma~\ref{lem:Normals} implies that $d \ge \frac{n}{2}$.  
% Thus $\frac{n}{2} \le d \le 6$ so 
% $n \le 12$.
% Furthermore:
%}

%%%%%%%%%%%%%%%%%%%%%%%%%%%%%%%%%Figure Begin
\begin{figure}[htbp]
\centering
\includegraphics[width=0.5\textwidth]{Figures/ProjectionPrivate}
%\includegraphics[width=0.7\textwidth]{Figures/ProjectionPrivate}
\caption{Projection of downward face $f$ lies within an equilateral triangle outside base edge $e$. Here $P_n=P_7$.
}
\figlab{ProjectionPrivate}
\end{figure}
%%%%%%%%%%%%%%%%%%%%%%%%%%%%%%%%%Figure End




Finally, we will rule out odd $n > 6$ by proving the following lemma.

%\JOR{Not sure why a Claim rather than a Lemma...?}
%\Anna{just that it's nested inside a lemma.  I don't mind calling it a lemma}
%\JOR{Oh, I see, but still I think better a lemma.}
\begin{lemma}
\label{lem:all-down}
%\Anna{oops, we need $\D$ in the lemma statement}
For any
dome over an equiangular $n$-gon with odd $n \ge 6$, % 
%all the dome faces incident to base edges have downward normals, i.e., 
$d \ge n$.
\end{lemma}


%Therefore there is no domeable equiangular $n$-gon with odd $n \ge 6$, since we would need $n \le d \le 6$.
%\end{proof}

%\anna{
%\begin{proof}[Proof outline for Lemma~\ref{lem:all-down}.]
%Each base vertex $b_i$ has four incident faces and thus, considered in isolation, has only one degree of freedom for the dihedral angles of incident edges. 
%Let $\d_i$ be the dihedral angle of $\D$ at base edge $e_i = b_i b_{i+1}$.  Because $b_i$ and $b_{i+1}$ share edge $e_i$ and have the same face angles, we can show that $\d_{i-1}=\d_{i+1}$.  
%Since $n$ is odd, this implies
%that all the $\d_i$'s are equal. 
%Lemma~\lemref{Normals} implies that at least one $\d_i$ is $> 90^\circ$, so they all are. 
%\end{proof}
%}
We begin with a preliminary result that will also be useful later on.
%
We henceforth often abbreviate ``dihedral angles'' by \defn{dihedrals}.
\begin{lemma}
\label{lem:dihedrals}
Let $P$ be an equiangular $n$-gon, $n \ge 6$ ($n$ may be odd or even) with edges $e_1, e_2, \ldots, e_n$.
Let $\D$ be a dome over $P$ and let $\d_i$ be the dihedral angle of $\D$ at base edge $e_i$.
If $n$ is even then the dihedrals $\d_i$ for $i$ even are all equal and the dihedrals $\d_i$ for $i$ odd are all equal.  If $n$ is odd then all dihedrals $\d_i$ are equal.
\end{lemma}
%
\begin{proof}
Let the vertices of $P$ be $b_i$, with $e_i = b_i b_{i+1}$. 
Since $P$ is equiangular and $n \ge 6$, the angle $\b$ at $b_i$ is $\ge 120^\circ$. 
By Lemma~\lemref{3Triangles}, there are three dome triangles $t_{i,1}, t_{i,2}, t_{i,3}$ incident to $b_i$, where $t_{i,1}$ is incident to $e_{i-1}$ and $t_{i,3}$ is incident to $e_i$;
see Fig.~\figref{Claim37}.  Each triangle has a $60^\circ$ angle at $b_i$. Together with the base angle of $\beta$, this gives---in isolation---one degree of freedom for the dihedral angles 
of edges incident to $b_i$, ranging from coplanarity of the first two triangles to coplanarity of the last two triangles.

Because there is one degree of freedom at $b_i$, 
fixing $\d_i$ then fixes the
dihedral angles of all edges 
incident to $b_i$ and to $b_{i+1}$. Furthermore, by reflective symmetry of $b_i$ and $b_{i+1}$ (which have the same face angles), 
we have $\d_{i-1}=\d_{i+1}$.
This implies the result.
%Thus even dihedrals are equal, odd dihedrals are equal, and for $n$ odd, all dihedrals are equal.
\end{proof}


\begin{proof} [Proof of Lemma~\ref{lem:all-down}]
By Lemma~\ref{lem:dihedrals}, all the dihedrals of $\D$ at base edges of $P$ are equal. 
By Lemma~\lemref{Normals}, at least one dome face incident to a base edge has a downward normal (equivalently, a dihedral angle $> 90^\circ$). Thus all the dome faces incident to base edges must
have downward normals.
\end{proof}

This completes the proof of Lemma~\ref{lem:12-and-6}.



%see Fig.~\figref{Claim37}.
%%%%%%%%%%%%%%%%%%%%%%%%%%%%%%%%%Figure Begin
\begin{figure}[htbp]
\centering
\includegraphics[width=0.60\textwidth]{Figures/Claim37}
\caption{Dome triangles incident to base vertex $b_i$.
}
\figlab{Claim37}
\end{figure}
%%%%%%%%%%%%%%%%%%%%%%%%%%%%%%%%%Figure End

\section{Proof of Theorem~\thmref{EquiAng}(b): Restrictions on edge lengths}
\label{sec:main-part-b}

We now turn to constraints on the edge lengths of equiangular polygons that can be domed,
part~(b) of our main theorem.
Of course all edge lengths must be 
integers
because the dome is composed of unit equilateral triangles.
For $n \le 6$, the constraints are relatively easy to identify. 
For $n=8,10,12$, significant argument is needed to establish if-and-only-if characterizations.
We first examine the edge length conditions in more detail.


\subsection{List of Conditions}
\seclab{ListConditions}
%
Recall the conditions stated in the main theorem:
for $n=3,4,5,6$ every equiangular integer polygon can be domed, and for $n=8,10,12$ an equiangular integer $n$-gon can be domed iff the lengths of the odd-numbered edges are equal and the even-numbered edge lengths are those of an equiangular $\frac{n}{2}$-gon.

Combining this with the known  characterizations of equiangular integer $n$-gons for $n=3,4,5,6$~\cite{ball2002equiangular}, 
we obtain the following:
%
\begin{description}
\item[$n=3$:] All side lengths equal.
\item[$n=4$:] Side lengths $a,b,a,b$, i.e., all 
integer rectangles. See Lemma~\lemref{Rect}.
\item[$n=5$:] All side lengths equal. This follows because every  
integer equiangular pentagon is regular~\cite{ball2002equiangular}.
%Here is an earlier paper that proves that an equiangular pentagon with integer side lengths is regular.
%Ball, Derek. “Equiangular polygons.” The Mathematical Gazette 86.507 (2002): 396-407.
\item[$n=6$:] Side lengths are
those of a polyiamond $6$-gon, namely, $a,b,c,a',b',c'$ with $a-a' = b'-b = c-c'$:
Fig.~\figref{Polyiamond6}.
%This is a polyiamond.
\item[$n=8$:] Odd edges have equal length $a$.
Even edges have lengths $b,c,b,c$. 
%$a+1 \le b,c$. 
% \Anna{I don't think we need $a+1 \le b,c$---see Coauthor ``An extra edge-length constraint?'' where we debated this.}
%\JOR{I looked at coauthor and thought through the $a+1$ issue, and I am now on board.
%I think I was mixing the base edge lengths with the
%edge lengths one story up.}
\item[$n=10$:] Odd edges have equal length $a$, even edges have equal length $b$. 
%$a+1 \le b,c$.
For example, see the earlier Fig.~\figref{IcosaSlicesDifferent}(b).
\item[$n=12$:] Odd edges have equal length $d$. 
% \Anna{Is it easy to switch to: odd edges have length $a$, to be consistent with the other cases? If this means lots of changes, let's not!}
% \JOR{I think I'd rather not, because $a,b,c$ is natural
% for those eqns and Fig.~\figref{Polyiamond6}.}
%
% \Anna{The following condition goes away: $d+1 \le a,b,c,a',b',c'$.}
Even edges are those of an equiangular $6$-gon, i.e., lengths are $a,b,c,a',b',c'$ with $a-a' = b'-b = c-c'$.
\end{description}

We aim to prove that an equiangular $n$-gon can be domed if and only if
the above constraints are satisfied.
The cases $n=3,4,5$ are already settled.

\subsubsection{$n=6$}
\seclab{n=6}

We prove the following more general result:

% \Anna{I've removed some stuff here and instead added to the proof.}
%\Anna{The point is that any polyiamond is made by truncating $0,1,2$ or $3$ vertices off an equilateral triangle.}
\hide{For $n=6$, the side-length constraints imply that $P_6$ is a polyiamond,
a truncation of the three corners of an equilateral triangle: Fig.~\figref{Polyiamond6}.
\begin{align*}
c' + a + b  = \; & b + c + a' \\
a -a'  = \; &  c - c' \\
a' + b' + c'  =\;  & b + c + a' \\
b' - b  =  \;& c - c'
\end{align*}
%see Fig.~\figref{Polyiamond6}.
} % end hide



\begin{lemma}
\lemlab{Polyiamond}
Every convex polyiamond can be domed.
\end{lemma}
\begin{proof}
A polyiamond on $3,4,5,6$ sides is formed by cutting off $0,1,2,3$, respectively, corners (small equilateral triangles) off a large integer equilateral triangle $T$; see Fig.~\figref{Polyiamond6}.
%\Anna{This does all $n$ at once:}
%\JOR{Nice!}
Construct a regular tetrahedron with triangle $T$ as one face.
Fig.~\figref{fig-polyiamonds} shows the case of a 6-sided polyiamond.
Extend the up-to-three small cut-off triangles to regular tetrahedra, shaded gray in the figure.  Cutting the gray tetrahedra off the large tetrahedron results in a truncated tetrahedon 
with the initial (pink) polyiamond as one face.
All the other faces are integer polyiamonds as well.
\end{proof}

%%%%%%%%%%%%%%%%%%%%%%%%%%%%%%%%%Figure Begin
\begin{figure}[htbp]
\centering
\includegraphics[width=0.4\textwidth]{Figures/Polyiamond6}
%\includegraphics[width=0.75\textwidth]{Figures/Polyiamond6}
\caption{A polyiamond in an integer equilateral triangle
of side length $a+b+c'$ (here any of $a', b, c'$ may be $0$). 
%$a -a' =  b' - b = c - c'$.
}
\figlab{Polyiamond6}
\end{figure}
%%%%%%%%%%%%%%%%%%%%%%%%%%%%%%%%%Figure End

%%%%%%%%%%%%%%%%%%%%%%%%%%%%%%%%%Figure Begin
\begin{figure}[htbp]
\centering
\includegraphics[width=0.4\textwidth]{Figures/fig-polyiamonds}
%\includegraphics[width=0.75\textwidth]{Figures/fig-polyiamonds}
\caption{The (pink) polyiamond is domed by the truncation of three corner tetrahedra.}
\figlab{fig-polyiamonds}
\end{figure}
%%%%%%%%%%%%%%%%%%%%%%%%%%%%%%%%%Figure End


\hide{
%\Anna{The proof should be fixed to handle $n=3,4,5$ polyiamonds too.}
%\JOR{Pepa's proof:
%``given a pink polyiamond, extend every other side to make an equilateral triangle, extend that triangle to a tetrahedron, and cut off 3 smaller tetrahedra at its corners to produce a solid that has the pink polyiamond as its face.''}
%\Anna{The following still needs fixing.}
We start with the case $n=6$
\anna{where the polyiamond $P_6$ is formed by cutting three corners off a large equilateral triangle.
Take the large triangle to be one face of a regular tetrahedron. 
See Fig.~\figref{fig-polyiamonds}.
Extend the three small cut-off triangles to regular tetrahedra, shaded gray in the figure.  Cutting the gray tetrahedra off the large tetrahedron results in a truncated tetrahedon 
with the (pink) polyiamond as one face.
All the other faces are integer polyiamonds as well.
}
%Draw the polyiamond $P_6$ on one face of a regular tetrahedron so that every other edge
%(even edges) shares boundary with that face.
%See Fig.~\figref{fig-polyiamonds}.
% Anna: Pepa's idea of "extending" the even sides was already accomplished by putting $P_6$ on one face of a tetrahedron.  So I've reworded a bit. 
%Extend the even edges to equilateral triangles in that face, and extend those equilateral triangles to regular tetrahedra,
%shaded gray in the figure.
%Cut off these three tetrahedra.
%This leaves the (pink) polyiamond a face of the truncated polyhedron. % $\P$.
%$\P$ only used once, so now removed.
%All the other faces of 
%$\P$ 
%the polyhedron
%are polyiamonds, and so composed of equilateral triangles.
Therefore, we have domed $P_6$.

\jor{To dome $P_n$ for $n=3,4,5$, 
follow the $P_6$ construction,
except do not remove (as appropriate)
either %$3$, $2$, or $1$
three, two, or one
of the gray tetrahedra in Fig.~\figref{fig-polyiamonds}.
Again the remaining faces are polyiamonds, and the
conclusion follows.}
\end{proof}
%see Fig.~\figref{fig-polyiamonds}.
} % end hide

\subsection{Conditions $\to$ domeable} %, for $n=8,10,12$}
For $n=8,10,12$, we first show that the edge-length conditions listed above
lead to domings of each $P_n$---the ``if'' direction of the claim.
The construction begins the same way in all three cases.
%All three are similar in the following sense.
Suppose the odd base edges have length $\ell$.  Erect an equilateral triangle $t$ with side length $\ell$ on each odd base edge.
% Anna: an $\ell times \ell$ triangle doesn't make sense.
%
%We erect an $\ell \times \ell$
%equilateral triangle $t$ over
%each odd base edge, which are all of length $\ell$.
This triangle has a downward normal.
Two consecutive triangles are joined by
%Join to either side of $t$ 
an upward facing trapezoid.
Then at the height of the apex of $t$,
we have a planar polygon on $n/2$ vertices---a ``lid.''
% an $n/2$-gon ``roof.''
% \Anna{Hmm. Isn't the ``roof'' the dome (or specifically the dome of a rectangle) rather than the polygon?  Maybe ``planar polygon on $n/2$ vertices.''}
For an octagon, 
the %roof 
lid is a rectangle, which we then dome.
For a decagon,
the %roof 
lid is a regular pentagon, which we then dome.
For a dodecagon,
the %roof 
lid is a polyiamond hexagon, which completes the dome. 
Below we give further details and show that the resulting polyhedron is convex in each case.

\subsubsection{$8$-gons $\to$ domeable}
We dome any octagon satisfying the edge-length constraints as illustrated
in Fig.~\figref{Octa_4x2}. The four downward pointing faces are tilted at
dihedral angle the same
as the corresponding faces of a square antiprism 
(the case that arises from a regular octagon).
One level of triangles results in a rectangular \anna{lid},
%boundary, 
which we then ``roof'' as 
seen in Figs.~\figref{RectInt} and~\figref{Rect_3x1_roof}.
%\Anna{We should say why the result is convex.}
% Anna: Great!  I didn't know this.
Combining the dihedral angles of a square antiprism
($\approx 104^\circ$) and a square pyramid
($\approx 55^\circ$)
leads to a dihedral angle where the roof meets a
side of $\approx 159^\circ < 180^\circ$ and so is convex.
%https://polytope.miraheze.org/wiki/Square_antiprism
%https://polytope.miraheze.org/wiki/Square_pyramid

%see Fig.~\figref{Octa_4x2}.
%%%%%%%%%%%%%%%%%%%%%%%%%%%%%%%%%Figure Begin
\begin{figure}[htbp]
\centering
\includegraphics[width=0.6\textwidth]{Figures/Octa_4x2}
%\includegraphics[width=0.75\textwidth]{Figures/Octa_4x2}
\caption{An octagon base with edge lengths $1,4,1,2,1,4,1,2$.
The 
%one-storey height 
``lid'' is a $4 \times 2$ rectangle.}
\figlab{Octa_4x2}
\end{figure}
%%%%%%%%%%%%%%%%%%%%%%%%%%%%%%%%%Figure End


\subsubsection{$10$-gons $\to$ domeable}
For certain values of $a,b$,
the length pattern $a,b,\ldots,a,b$ can be achieved by adjusting the height of the
slice through an icosahedron; see
Figs.~\figref{SliceFigs}(c) and~\figref{IcosaSlicesDifferent}(b).
But the more general strategy was described above:
% Anna: we don't need a \le b
%If $a \le b$, 
At the height of the $a$-triangle's apex,
the lid is
%we have 
a regular pentagon with side length $b$.
This can then be domed 
with a pentagonal pyramid
as in Fig.~\figref{Pyramids345}.
%\Anna{We should say why the result is convex.}
Combining the dihedral angles of a pentagonal antiprism
($\approx 101^\circ$) and a pentagonal pyramid
($\approx 37^\circ$)
leads to a dihedral angle where the top 
meets a
side of $\approx 138^\circ < 180^\circ$ and so is convex.
%https://polytope.miraheze.org/wiki/Pentagonal_antiprism
%https://polytope.miraheze.org/wiki/Pentagonal_pyramid

\subsubsection{$12$-gons $\to$ domeable}
The odd edges of equal length $d$ are covered by downward facing equilateral $d$-triangles.
See Fig.~\figref{PepaConstruction}.
One level up, 
%leaves 
the lid is
%the even edges 
%form
a polyiamond $6$-gon, 
which can be filled in to create
%creating 
the last face of the dome.
%which
%follows the patterns previously shown in Fig.~\figref{Polyiamond6}.
%See Fig.~\figref{PepaConstruction}.
%%%%%%%%%%%%%%%%%%%%%%%%%%%%%%%%%Figure Begin
\begin{figure}[htbp]
\centering
\includegraphics[width=0.5\textwidth]{Figures/PepaConstruction}
%\includegraphics[width=0.75\textwidth]{Figures/PepaConstruction}
\caption{Equiangular $12$-gon.
Red hexagon is polyiamond roof.
Green polygon is base $P_{12}$, angle $\b = 150^\circ$.
$a,b,c=1,2,1$; $d=1$.}
\figlab{PepaConstruction}
\end{figure}
%%%%%%%%%%%%%%%%%%%%%%%%%%%%%%%%%Figure End



%\clearpage
\subsection{Domeable $\to$ conditions}

In this section we prove
that any domeable equiangular integer $n$-gon, $n=8,10,12$, satisfies the edge length conditions. We know from Lemma~\ref{lem:6dome} that
a dome $\D$ over such an $n$-gon has
%if $P$ can be domed then it has a dome $\D$ with %no $V_3$ vertex and thus with 
at most 6 dome vertices. We will prove properties of 
the 
%such a dome $\D$ 
dome to derive the edge length conditions.

%\Anna{We need to explain that the argument is: if $P$ can be domed then it has a dome $\D$ with no $V_3$ vertex, and we prove properties of THAT dome $\D$ to get the edge length conditions on $P$.}
%\JOR{I wish I had a collection of examples that have $V_3$
%vertices. There seem to be only a few? E.g., square base, two $V_3$
%vertices, whose removal leads to a square pyramid.}
%\Anna{Simplest is a tetrahedron with an extra tetrahedon glued to one face.  Or what was the one you asked me about the other day?}
%\Anna{We do not at all characterize ALL domes -- that's an open question for later.  Only all equiangular bases.}


\subsubsection{$12$-gons $\to$ conditions}
We first handle the $n=12$ case, as it is simpler than $n=8,10$.
We will repeat some of the logic when addressing $n=8,10$ below.
We present the argument in five steps.  Refer throughout to Fig.~\figref{n12Proof}.

%\Anna{I notice that starting here, we have reversed odd and even base edges. In the main theorem and everthing up to this point, the odd edge lengths are equal (and they have the downward faces).  }
%\Anna{OK, I think it's all fixed.}
\begin{enumerate}[(1)]
%
\item
As we saw in Lemma~\ref{lem:dihedrals},
the dihedrals at even base edges are equal, and the dihedrals at odd base edges are equal.  By Lemma~\ref{lem:Normals}, one set, say the 
odd
%even 
edges, 
has dihedrals $> 90^\circ$, i.e., the corresponding dome faces have downward normals.
%\Anna{to be precise it's not the dihedrals that point downward}
%As we saw in \jor{Lemmas~\ref{lem:dihedrals}
%and~\ref{lem:all-down},}
%the dihedrals for even edges are equal, and they have downward normals 
%and the dihedrals for odd edges are equal.
%See Fig.~\figref{n12Proof}.
%%%%%%%%%%%%%%%%%%%%%%%%%%%%%%%%%Figure Begin
\begin{figure}[htbp]
\centering
\includegraphics[width=0.65\textwidth]{Figures/n12Proof}
%\includegraphics[width=1.0\textwidth]{Figures/n12Proof}
\caption{Faces on odd
%even 
edges are triangles pointing down.
Faces on even %odd 
edges are quadrilaterals.
% \Anna{the $120^\circ$ is for the WHOLE angle (the union of two $60^\circ$ angles).  Could you please add little circular arcs to indicate this? The arrow makes it seem like the wrong thing.}}
}
\figlab{n12Proof}
\end{figure}
%%%%%%%%%%%%%%%%%%%%%%%%%%%%%%%%%Figure End
%
\item The projection argument in %Section~\secref{NormalsPrivate}
Lemma~\ref{lem:d-le-6}
shows that 
the $6$ faces incident to odd base edges have disjoint sets of dome vertices.
%each face on an odd 
%edge needs at least one %private vertex.
Since 
our dome $\D$ has
at most $6$ dome vertices, 
%by Lemma~\ref{lem:6dome},
each of these faces gets one 
%private vertex, 
dome vertex, and so it must be a triangle.
Also, each dome vertex is the apex of one of these triangles.
%\Anna{This is a stronger claim than what we originally wrote: we need that ANY dome vertex of a face on an odd edge must be private.}
%
\item Each base vertex \anna{$b$} has three incident dome triangles.  One of them is in the downward-normal face at the vertex.  The edge incident to $b$ that lies between the other two dome triangles will be called a ``red'' edge (dashed in Fig.~\figref{n12Proof}).
Next we prove that the 12 red edges 
%in Fig.~\figref{n12Proof} 
are flat, i.e., have a $180^\circ$ dihedral.
%\Anna{I put the 12 up front since it's independent of nonflat.}
By symmetry
of the base vertices, all the red edges have the same dihedral, so if one is nonflat, then they all are.  
%Then we have $12$ red edges.  
What are the 
dome vertices at the other ends of these nonflat red edges? (Note that a red edge cannot go to another base vertex.)
%\Anna{just hoping to clarify that we follow the edge past flat dome vertices}
%other endpoints of these red edges?  
%They cannot be base vertices, so they must be dome vertices. 
Since there are $12$ red edges and only $6$ dome vertices,
some dome vertex $u$ is an endpoint of at least two red edges. 
%Because $u$ is a private vertex,
As noted above, $u$ is the apex of a downward-normal triangle, so
it 
also has two (different) edges to the endpoints of an odd
%even 
base edge. Because $u$ has at most $5$ incident dome triangles, 
it has at most one more incident dome edge
beyond the two red edges and two edges to the base.
%\JOR{Reviewer~4 thinks these sentences are confusing. Repaired now.}
This means that two of $u$'s edges separated by $60^\circ$
(thus in one face)
are incident to non-consecutive base vertices. 
But it is impossible for two non-consecutive base vertices to lie in the same dome face.
%But this is impossible,
%as it would force collinearities in base edges. 
Thus the red edges cannot end at dome vertices.
Thus they must be flat edges.
The face on any 
even %odd 
base edge then has $120^\circ$ face angles at the base vertices.  We next argue that each such face is a quadrilateral.
%
\item Consider now a dome vertex $v$, the apex of a triangle face on an odd
%even 
edge, say $e_1$ in the figure.
Two of its (at most) five incident edges go to the base endpoints of 
$e_1$. 
%$e_2$.
And the two 
incident edges to the
left and right, $e_l$ and $e_r$,
(whether flat or not) must be horizontal and parallel to the corresponding base edge
$e_0$ and $e_2$.
%($e_1$ and $e_3$). \Anna{Fix!}
The angle between $e_l$ and $e_r$ is then $120^\circ$.
Therefore $e_c$, the central edge between them, must be flat, and $e_l$ and $e_r$ are not flat.
%
\item 
The face on any even
%odd 
base edge $e_i$ is then a trapezoid %quadrilateral 
with $120^\circ$ angles at the base and $60^\circ$ angles at the dome vertices.   
Then the dome vertices must all be at the same level, 
so the triangles on odd base edges all have the same edge length, say $\ell$. 
The length of the top edge of the trapezoid on $e_i$ is $|e_i| + \ell$.
These top edges 
joining the dome vertices then form an equiangular integer hexagon, a polyiamond, whose edge lengths must then be $a,b,c,a',b',c'$ with $a-a' = b'-b=c-c'$.  
%These top edges 
%form an equiangular hexagon joining the dome vertices, whose
%edge lengths satisfy the constraint  stated in the 
%claim above.  
Thus the 
even  
base edge lengths also satisfy that constraint, which proves the result for $n=12$.
%
\end{enumerate}

\subsubsection{$8,10$-gons $\to$ conditions}
\seclab{n810Proof}
Compared to the case $n=12$ there are some hurdles to overcome. 
One is that a 
downward-normal face on an odd  edge 
might
not be a triangle---it %may 
could
be a quadrilateral with two
dome vertices.
%private vertices.  
Another is that there may be 
dome vertices that are not included in any of the faces on the odd edges,
%non-private dome vertices, 
at most one for $n=10$ and at most two for $n=8$.  
Despite these difficulties, we will still 
follow the same outline as above, first proving
that the red edges in Fig.~\figref{n12Proof} are flat, and then analyzing the situation at one dome vertex.

%
\hide{These two cases are more difficult than $n=12$, because we can no longer argue (in step (2) above)
\Anna{Should we say: that the faces on even edges are triangles?  I mean, aren't we stuck even if all dome vertices are private?}
\JOR{Not sure how to respond to this. The original seems
to me to be true, and reads well.
This is just an introduction, so we don't need
absolute precision. As long as the steps to follow
are correct.}
% Anna: the original reads better!
%\jor{we no longer can}
% JOR: OK, I was trying to avoid splitting the verb,
% but sound is important.
%conclude 
that every dome vertex is private,
which was crucial for proving that red edges are flat.
For $n=10$ we may have one non-private dome vertex.  For $n=8$ we could have two non-private dome vertices.
} % end hide



\begin{enumerate}[(1)]
%
\item
As we saw in Lemma~\ref{lem:dihedrals},
the even %odd 
base edges have equal dihedrals, and the odd %even 
base edges have equal dihedrals and are incident to faces with downward normals.
%
%\item 
These downward faces on the odd %even 
edges, which we call %``blue''
\emph{blue}
faces, have $60^\circ$ angles at the base
by Observation~\ref{obs:down-face}. 
Thus each blue face is
either a triangle
with one 
dome
%private 
vertex, 
or a trapezoid with two 
dome
%private 
vertices (again, following the projection argument in 
Lemma~\ref{lem:d-le-6}).
%Section~\secref{NormalsPrivate}).
See Fig.~\figref{RedFaces_FigB}.
%\anna{We call the dome vertices of blue faces the \emph{blue} vertices.}
For $n=12,10,8$, we need at least $n/2=6,5,4$
of the $6$ dome vertices to lie in blue faces.
%private vertices
%out of the $6$ dome vertices.
For $n=12$, 
this means that the 
blue
%downward 
faces are all triangles.
For $n=10,8$, there could be one or two
blue faces that are trapezoids.
%downward quadrilateral faces.
%
% \item 
% \Anna{I removed the observation, as I explain below -- I do not see a way make (long) flat edges well-defined.  We should restrict our use of ``nonflat'' to edges between dome triangles.  Those make sense.  In the context where we need the observation (below), things are precise.}
% \JOR{Good point that long nonflat edges are difficult
% to define. And now the argument doesn't need the.}
%\label{item:consecutive-edge-obs}
% \Anna{Please see below.  I want to remove the Observation from here and just make the argument in-place below.  The reason is that ``non-flat'' edges are not a well-defined thing.  (Nothing is flawed, it's just a question of writing it clearly.)}
%
%We again need the observation that a dome vertex cannot have two
%incident edges (flat or nonflat) that are separated by $60^\circ$ (thus in one face) and that go to base vertices (or points on the base) that do not lie in the same base edge.
%\Anna{This is a bit subtle and at least one of my re-wordings was wrong.}
%\JOR{Seems clear now.}

%consecutive (separated by  $60^\circ$) edges (flat or nonflat)
%that go to base points 
%that are not consecutive.
% Anna: to keep the wording the same as for n=12 case.
%not both on the same edge.
%
%\Anna{hmm, do we need the possibility of flat? We didn't in the $n=12$ case.  I'll think about this.}
%\JOR{If the edges of $v$ go to base points not both on the same edge, then between those base points there must
%be a base vertex $b$. Then there is an edge from $b$
%to $v$ splitting $60^\circ$. Contradiction. I am
%not seeing how flat/nonflat makes a difference.}
%\Anna{The observation is true for flat or nonflat edges.  I was wondering if we needed the case of flat edges (we didn't for $n=12$). The answer is Yes, we do. }
%
\item 
%Using this observation, 
As for the $n=12$ case above, we identify a ``red'' edge at each base vertex $b$---the edge between the two dome triangles at $b$ that are not in the blue face. As before, we will argue that all red edges are flat. 
%The next step is to prove that all the red edges
%in Fig.~\figref{n12Proof} are flat. 
Suppose not.  As above, by symmetry, this implies that all the red edges are nonflat.
Here is where matters are more complicated
than in the $n=12$ case.
Consider the faces incident to the even %odd 
base edges. 
Because the red edges are nonflat, each such face is bounded by red edges at the base and we call it a \emph{red} face.  A red face has  
$60^\circ$ degree face angles at the base, so it must be a triangle with one dome vertex, or a trapezoid with two dome vertices.
% \Anna{I omitted the definition of red vertices.}
%We call the dome vertices of red faces the \emph{red} vertices.
%\anna{We have not yet ruled out a vertex being both blue and red.}
See Fig.~\figref{RedFaces_FigB}.
Note that there is an upward-normal face between consecutive red and blue faces.

%see Fig.~\figref{RedFaces_FigB}.
%%%%%%%%%%%%%%%%%%%%%%%%%%%%%%%%%Figure Begin
\begin{figure}[htbp]
\centering
\includegraphics[width=0.55\textwidth]{Figures/RedFaces_FigB}
\caption{
If red edges are nonflat, then faces incident to the base alternate between red faces and blue 
downward faces, each with one or two dome vertices.
}
\figlab{RedFaces_FigB}
\end{figure}
%%%%%%%%%%%%%%%%%%%%%%%%%%%%%%%%%Figure End

Our next goal is to show that 
a dome vertex cannot lie in more than one red or blue face.
%all the private and red vertices are distinct.  
This immediately gives a contradiction, since we need at least 
one dome vertex for each of the at least $8$ base edges
%$8$  private/red 
%vertices (one per base edge) 
but there are at most $6$ dome vertices.  

\begin{claim}
A dome vertex cannot be shared by two red/blue faces.
\end{claim}

\begin{proof} Suppose dome vertex $v$
is shared by %two red/blue faces, 
two faces, either of which may be red or blue,
say on base edges $e_i$ and $e_j$.
See Fig.~\figref{n810RedFace}.
In case either of these faces is a trapezoid, we focus on the equilateral triangle (a union of dome triangles) that is contained in the trapezoid and  has its top $60^\circ$ corner at $v$ and its other two ``base'' corners on the base edge.  
We consider what happens between the two triangles.  If they share a side, then the common base corner is a
base vertex with only two incident dome triangles, a contradiction.
Next, suppose 
there is a $60^\circ$ angle at $v$ between the two triangles. Then this angle is part of one dome face  that extends to a base corner of each of the two triangles. This is only possible if the two base points lie on a common base edge joining $e_i$ and $e_j$.  But then we have base vertices with only two incident dome triangles, 
again a contradiction.   
Thus there must be at least $120^\circ$ of surface angle at $v$ on each side between the two triangles, a contradiction to $v$ having at most $5$ incident dome triangles.
\end{proof}
%
%see Fig.~\figref{n810RedFace}.
%%%%%%%%%%%%%%%%%%%%%%%%%%%%%%%%%Figure Begin
\begin{figure}[htbp]
\centering
\includegraphics[width=0.75\textwidth]{Figures/n810RedFace}
\caption{(a)~A dome vertex $v$ cannot be in a red face and a blue face.
%A red vertex cannot coincide with a private vertex $v$. 
(b)~Two red faces cannot share a vertex $v$.
% \Anna{If my condensed proof above makes sense (does it?), then please remove the little blue triangle on $e$ in part (b). There is not too much reason for having (a) and (b) now, except that there is the parity issue for (a) ($j$ even, $i$ odd implies an even number of base edges between them (which turns out not to be needed in the proof after all)).  Still, I suggest leaving (a) and (b).}
%\JOR{I will study the proof...}
}
\figlab{n810RedFace}
\end{figure}
%%%%%%%%%%%%%%%%%%%%%%%%%%%%%%%%%Figure End


\hide{
Our next goal is to show that there are not enough dome vertices to provide these red vertices. 
%
%\Anna{I find Figure B here \url{https://coauthor.csail.mit.edu/file/vzDXismkwzGbvmWy7} very helpful if you have time and energy for it.}
% \JOR{Followed FigB even different sizes of
% triangles and traps.}
% \Anna{Great!}
%Call each face bound by red edges at the base \emph{red faces}.
%Each has one or two red vertices.
We first prove (in (a) below) that a red vertex cannot be a private vertex.  
For $n=12$ this %does it, 
settles matters, since all $6$ dome vertices are private, so there cannot be any red vertices.  For $n=10$ we need $5$ private vertices; there may be one non-private dome vertex left over, but we cannot have all $5$ red faces (with $10$ red edges) meeting at one vertex.  For $n=8$ we may have two non-private dome vertices and we cannot rule out shared red vertices just by counting. So we need a second claim~((b) below) to complete the argument.
%Now we address two claims.
\begin{enumerate}[(a)]
\item A red vertex cannot
\anna{be} 
%coincide with 
a private vertex.
See Fig.~\figref{n810RedFace}(a).

%\Anna{I've gone back to the longer version from coauthor}
Suppose dome vertex $v$ is a \anna{private vertex of the blue} face on base edge $e_j$, $j$ even, and also a red vertex of the red face on some edge $e_i$, $i$ odd. 
\anna{In case either of these faces is a trapezoid, we focus on the equilateral triangle (a union of dome triangles) that is contained in the trapezoid and  has its top $60^\circ$ corner at $v$ and its other two (``base'') corners on the base edge.  
We then have one red and one blue triangle incident to $v$; see 
Fig.~\figref{n810RedFace}(a)).
We consider what happens between the red and blue triangle.  If they share a side, then the common base corner is a
base vertex with only two dome triangles, a contradiction. 
If there is a $60^\circ$ angle at $v$ between the two triangles, then TO BE FILLED IN, a contradiction.  Thus there must be at least $120^\circ$ of surface angle at $v$ on each side between the red and blue triangles, a contradiction to $v$ having at most $5$ incident dome triangles.}

\Anna{I'm rewriting because I'm bothered by ``a non-flat edge'' from $v$ to the base -- in what sense is that an edge? It's just a straight segment composed of dome triangle edges.}

If base edges $e_i$ and $e_j$ were consecutive on the base, then the faces would share a base vertex $u$ and share vertex $v$ which means they must share the edge $uv$---\anna{a contradiction to $u$ having three incident dome triangles.}
%but
%they have different edges at $u$, a contradiction.  
Thus $e_i$ and $e_j$ are not consecutive.  This implies that $v$ has two edges (one of which may be flat) to points on $e_j$ and two different edges (one of which may be flat) to points on $e_i$.    Between these two pairs of edges there are two gaps (in the cyclic order of edges around $v$. Because $v$ has at most 5 incident triangles, one of the gaps must be a $60^\circ$ face angle at $v$ bounded by one edge to a point in $e_i$ and one edge to a point in $e_j$. 
\hide{The edges $e_i$ and $e_j$ in the figure cannot be consecutive, for then they would share
a red/blue edge, but there are no such edges.
Because $v$ has at most %degree-$5$, 
\jor{$5$ incident triangles,}
one of the gaps between its \jor{incident} edges must be $60^\circ$.
} % end hide
But this is impossible by %the
observation~\ref{item:consecutive-edge-obs} above.
%

\item Two red faces cannot share a vertex $v$.
See Fig.~\figref{n810RedFace}(b).
Again because $v$ %has 
\jor{is} at most degree-$5$, one of the gaps between its edges must be $60^\circ$.
%The observation 
\jor{Observation~\ref{item:consecutive-edge-obs}} implies that the two edges bounded by this $60^\circ$ must go to the same
base edge $e$. But then those edges encompass a private vertex $u$. 
The only possible third edge incident to $u$ is to $v$, but that would split the $60^\circ$ angle there,
a contradiction.
\end{enumerate}
} % end long hide

This completes the proof that the red edges are flat.
We conclude that
faces on even %odd 
base edges have
$120^\circ$ angles at the base, 
and upward normals, as shown
in Fig.~\figref{n12Proof} (except that the blue faces need not be triangles).
%
\item 
The next goal is to prove that every blue face is a triangle and that the other faces on base edges are trapezoids.  
We accomplish this by analyzing
the local structure at a %private vertex $v$, 
dome vertex $v$ in a blue face.
%which must 
%be a $V_4$ or $V_5$ vertex (i.e., with
%$4$ or $5$ incident dome triangles).
% \Anna{I've omitted that $v$ is a $V_4$ or $V_5$ vertex. 
% Once upon a time we excluded $V_3$ vertices, but now we don't.  Our proof handles them ok.}
%


Vertex $v$ lies in a blue downward face $f$,
say on base edge 
$b_1 b_2 = e_1$.
%as in Fig.~\figref{PrivateVertex}.
Suppose, without loss of generality, that 
$v b_1$
%$v b_2$ 
is an edge of face $f$; then segment 
$v b_2$
%$v b_3$ 
lies in $f$ but is an edge of $f$  only if $f$ is a triangle.
Name the edges (flat or nonflat) between the dome triangles at $v$ as shown in Fig.~\figref{PrivateVertex}(a):
in the 
clockwise
direction from 
$v b_1$
%$v b_2$ 
is the edge $e_l$; in the counterclockwise
direction, there is an edge in $f$, then the edge $e_r$; and if there is a fifth edge, it is $e_c$ between $e_r$ and $e_l$.  
% \Anna{a bit awkward!} 
The upward face sharing edge 
$v b_1$
%$v b_2$ 
has face angle $120^\circ$ incident to 
$b_1$.
%$b_2$.  
Edge $e_l$ is in or on the boundary of this face (we do not assume the face is a quadrilateral), and
thus $e_l$ is horizontal and parallel to 
$e_0$.
%$e_1$.
%
%see Fig.~\figref{PrivateVertex}.
%%%%%%%%%%%%%%%%%%%%%%%%%%%%%%%%%Figure Begin
\begin{figure}[htbp]
\centering
\includegraphics[width=1.0\textwidth]{Figures/PrivateVertex}
\caption{Dome vertex $v$ in a blue face $f$. 
%Private vertex $v$. 
Face $f$ cannot be a quadrilateral.
}
\figlab{PrivateVertex}
\end{figure}
%%%%%%%%%%%%%%%%%%%%%%%%%%%%%%%%%Figure End
\item We next rule out 
face $f$ being a trapezoid.
%the face to the other side of $f$ being a quadrilateral.
Suppose it is. 
Then the angle of $f$ at $v$ is $120^\circ$ and edge $e_r$ is an edge of trapezoid $f$, thus horizontal and parallel to $e_1$.
See Fig.~\figref{PrivateVertex}(b).
But this means that the angle between $e_l$ and $e_r$ is the same as the base angle $\b$
which is greater than $120^\circ$.
Even if $v$ 
is %degree-$5$, 
a $V_5$ vertex,
there is only one edge $e_c$ between $e_l$ and $e_r$,
not enough to cover $\b$.
Therefore $f$ is a triangle
%with edge $v b_2$, and the 
$b_1 b_2 v$, and
edge $e_r$ lies outside $f$.
%
\item 
The same argument as above for $e_l$ then shows that
$e_r$ is  horizontal and parallel to 
\anna{$e_2$}.
%$e_3$.
The angle between $e_l$ and $e_r$ is then $90^\circ$, $108^\circ$, or $120^\circ$ for $n=8,10,12$ respectively.  This angle is too large to be covered by a single triangle,
so it requires at least two triangles.
Therefore $v$ must be a $V_5$ vertex,
and the edge $e_c$ is uniquely determined 
and the edges $e_l$ and $e_r$ are nonflat.
%
\item 
To wrap up,
all the blue downward %base 
faces are triangles, and
their dome vertices
%the private vertices 
are $V_5$ vertices.
The faces on the even %odd 
base edges are 
%quadrilaterals 
trapezoids with face angles of $120^\circ$ at the base and angles of $60^\circ$ at the dome vertices. 
(Recall the earlier $n=10$ example in Fig.~\figref{IcosaSlicesDifferent}(b).)
The trapezoids have the same dihedral angle and join the 
%private vertices 
dome vertices of the blue face, so
all 
those dome vertices
%the private vertices 
must be at the same height from the base 
%which implies,
and
the odd %even 
base edges must have equal length, say $\ell$.
%\Anna{Is this what we called the length?}
%\JOR{Yes.}
At the level of 
%the private vertices, 
those dome vertices
we have an equiangular $4$-, $5$- or $6$-gon. 
The length of the edge parallel to base edge $e_i$ is
$|e_i| + \ell$.
%$|e_i| +1$.
%
\end{enumerate}
This completes the analysis of the edge-length constraints for $n=8,10,12$,
and so proves part (b) of Theorem~\thmref{EquiAng}.

\subsection{Not all domes are slices of general deltahedra}
\seclab{NotSlice}
%\JOR{Anna's argument nearly verbatim.}
From Section~\secref{n810Proof} above, we know that for $n=8$,
the even edges are those of a rectangle, $b,c,b,c$.
Furthermore, the construction of the dome is unique---in particular,
the faces on the odd edges must have downward normals and the faces on the even edges must have upward normals.
See Fig.~\figref{Octa_4x2}.

We now argue that the octagon domes where $b \neq c$ do not arise from slicing a generalized convex deltahedron.
If they did, it would be possible to construct domes on the top and the bottom of the octagon such that the two halves form a convex polyhedron. But on the bottom, by convexity, the faces on the odd edges would then have upward normals (after we flip it). Then the even edges would have to be the ones with downward normals so their lengths would need to be equal.
So the construction is not realizable if $b \neq c$.

A similar argument can be made for $n=12$.


%\clearpage
\section{Further Results and Open Questions}
\label{sec:further-results}
% Anna: this sentence seems redundant:
%We have proven our main result, 
%Theorem~\thmref{EquiAng}.



%-------------------------------------------------------------------------------------------------

%\section{Open Problems}
%\noindent
We have characterized equiangular domeable polygons.
When we drop the equiangular condition,
many open problems remain. Here are a few.
\begin{enumerate}[(1)]
%
\item Is there any convex $n$-gon with $n > 12$ that can be domed?
A rough bound is $n\le 55$: every dome vertex has curvature at least $\frac{\pi}{3}$ so there are at most 11 dome vertices; every base vertex must be adjacent to a dome vertex and dome vertices have degree at most~$5$.
We can prove the stronger bound $n \le 24$.
%\JOR{Leaving this as a claim.}
%
\item Is there any convex $7$-gon that can be domed?
Fig.~\figref{9-gon_3Views} shows a ``one-story'' doming of a 
(non-equiangular) $9$-gon.
An $11$-gon can be constructed similarly, 
but the same construction scheme fails for a $7$-gon.
%\JOR{Reviewer~4 wants to see the $11$-gon. I think not worth the effort.}
We do not know if there is a domeable $7$-gon via some other construction.
%We have constructed $9$- and $11$-gons 
%(non-equiangular)
%that can be domed.
%See Fig.~\figref{9-gon_3Views}.
%
\item 
Is there any non-equilateral triangle that can be domed?
Glazarin and Pak conjectured~\cite{glazyrin2022domes} that,
even under their looser conditions,  
an isosceles triangle with edge lengths $2,2,1$
cannot be spanned.
%
\end{enumerate}
%%%%%%%%%%%%%%%%%%%%%%%%%%%%%%%%%Figure Begin
\begin{figure}[htbp]
\centering
\includegraphics[width=0.8\textwidth]{Figures/9-gon_3Views}
\caption{The top (red) is a polyiamond of $13$ equilateral triangles.
The base (blue) is a $9$-gon with base angles $120^\circ$ and $150^\circ$.
}
\figlab{9-gon_3Views}
\end{figure}
%%%%%%%%%%%%%%%%%%%%%%%%%%%%%%%%%Figure End

\subparagraph*{Acknowledgments.}
We benefited from suggestions from six anonymous reviewers
and from comments at the EuroCG presentation~\cite{DomesEuroCG}.

%------------------------------------------------------------
%\bibliographystyle{alpha}
\bibliography{refs}
%------------------------------------------------------------
%-------------------------------------------------------------------------------------------------
%-------------------------------------------------------------------------------------------------
%-------------------------------------------------------------------------------------------------
%-------------------------------------------------------------------------------------------------
%-------------------------------------------------------------------------------------------------


\end{document}
%%%%%%%%%%%%%%%%%%%%%%%%%%%%%%%%%%%%%%%
%%%%%%%%%%%%%%%%%%%%%%%%%%%%%%%%%%%%%%%
%%%%%%%%%%%%%%%%%%%%%%%%%%%%%%%%%%%%%%%
%%%%%%%%%%%%%%%%%%%%%%%%%%%%%%%%%%%%%%%
%%%%%%%%%%%%%%%%%%%%%%%%%%%%%%%%%%%%%%%
%%%%%%%%%%%%%%%%%%%%%%%%%%%%%%%%%%%%%%%
%%%%%%%%%%%%%%%%%%%%%%%%%%%%%%%%%%%%%%%
%%%%%%%%%%%%%%%%%%%%%%%%%%%%%%%%%%%%%%%
%%%%%%%%%%%%%%%%%%%%%%%%%%%%%%%%%%%%%%%

see Fig.~\figref{xxx}.
%%%%%%%%%%%%%%%%%%%%%%%%%%%%%%%%%Figure Begin
\begin{figure}[htbp]
\centering
\includegraphics[width=0.75\textwidth]{Figures/xxx}
\caption{caption.}
\figlab{xxx}
\end{figure}
%%%%%%%%%%%%%%%%%%%%%%%%%%%%%%%%%Figure End


\noindent

\begin{enumerate}
\squeezelist
\end{enumerate}


\begin{lemma}
\lemlab{zzz}
\end{lemma}
\begin{proof}
\end{proof}

\begin{definition}
\end{definition}

\begin{theorem}
\thmlab{zzz}
\end{theorem}
\begin{proof}
\end{proof}

